\chapter{Tutorials}
\label{chap:tutorials}

This chapter walks through the worked examples that ship with \lupnt{}. Every
example exists in two forms: a Python notebook under \file{python/examples/},
which renders the results as 3-D scenes, time histories, maps and histograms,
and --- for most of them --- a standalone C++ program under
\file{cpp/examples/tutorials/} that computes the same quantities and prints them
as console statistics. The two are deliberately kept numerically equivalent: the
notebook is the exploratory surface, the \code{.cc} file is the
compile-time-checked twin that exercises the same library entry points from
C++.

The chapter is organised as one section per example. Each section states the
scenario the example configures --- epochs, orbital elements, constellations,
station sets, sensor noise --- with the actual numbers taken from the notebook;
walks through the computation with excerpts of the real source; then describes
\emph{what the example produces}: every figure it draws, what is on each axis,
the magnitudes involved, and the conclusion the reader is meant to draw. Each
section closes with a pointer to the Part~\ref{part:math} chapter that specifies
the underlying model. A consolidated index is given in
\cref{tab:tut-examples}, and \cref{sec:tut-publications} maps the examples onto
the publications whose methodology they implement.

\begin{lupntnote}
The examples are not unit tests. They are allowed to download data, to take
minutes to run, and to print approximate numbers. Every numeric result quoted in
this chapter is the value stored in the committed notebook output for one
particular run on one particular machine; a re-run with a different random seed,
a different SP3 vintage, or a different Earthdata snapshot will move the last
digits. The authoritative statements about accuracy live in the
Part~\ref{part:math} chapters and in the validation material of
Part~\ref{part:verification}.
\end{lupntnote}

% ===========================================================================
\section{Running the Examples}
\label{sec:tut-running}

Both flavours are driven by the \code{pixi} workspace defined in
\file{pixi.toml} at the repository root. The workspace's activation block sets
the environment variables the examples assume --- \code{LUPNT\_DATA\_PATH}
pointing at \file{data/LuPNT_data}, \code{PYTHONPATH} pointing at
\file{python/}, \code{PECSIMPY\_BASE\_PATH} at the plasma data, and
\code{LUPNT\_OUTPUT\_PATH} at \file{output/} --- so no manual export is needed
as long as commands are run through \code{pixi run}.

\subsection{Python notebooks}
\label{sec:tut-running-python}

The notebooks import \pylupnt{}, the pybind11 binding layer, which must be built
once. Building the bindings and registering a Jupyter kernel that inherits the
workspace environment is a single task each:

\begin{lstlisting}[language=shell, caption={Building and running the Python examples.}, label={lst:tut-run-python}]
# from the repository root
pixi run build-py         # builds the pylupnt-dev extension module
pixi run install-kernel   # registers a Jupyter kernel with LUPNT_DATA_PATH baked in

# then open any notebook with the "LuPNT (pixi)" kernel, e.g.
#   python/examples/ex1_propagate_orbit.ipynb

# or execute them headlessly and get a pass/fail summary
pixi run run-notebooks -- --only ex1 --timeout 600
pixi run run-notebooks                      # all of ex1 .. ex17
\end{lstlisting}

Every notebook begins with the same defensive preamble: it walks up from the
working directory until it finds \file{python/pylupnt/__init__.py}, puts that
\file{python/} first on \code{sys.path}, and pins \code{LUPNT\_DATA\_PATH} to the
sibling \file{data/LuPNT_data}. This makes a notebook use \emph{this} checkout's
bindings and data even when a stale copy sits on a global \code{PYTHONPATH} or
in an old kernelspec.

\code{run-notebooks} executes each notebook end-to-end without writing outputs
back. The companion task \code{render-notebooks} executes them \emph{in place},
so that the Sphinx build (which never re-executes notebooks) shows real results;
\code{render-tutorials} additionally optimises the embedded figures. Several of
the heavier examples (\cref{sec:tut-ex6}, \cref{sec:tut-ex7},
\cref{sec:tut-ex8}, \cref{sec:tut-ex9}, \cref{sec:tut-ex10},
\cref{sec:tut-ex11}) delegate the expensive run to a sibling
\code{ex*\_run\_*.py} script and let the notebook read back the exported results,
which keeps the rendered tutorial fast. Those scripts seed their output
directory from a shipped copy under \file{data/LuPNT_data/examples/}, so the
notebook still plots real numbers on a machine that has never run the
simulation.

\subsection{C++ tutorial binaries}
\label{sec:tut-running-cpp}

The directory \file{cpp/examples/} is its own CMake project, \code{lupnt\_examples}.
Its \file{CMakeLists.txt} globs every \code{.cc} file under \file{tutorials/},
turns each into an executable named after the file stem, links it against
\code{lupnt::lupnt} plus TBB, and collects them all under a convenience target
\code{all\_examples}:

\begin{lstlisting}[language=shell, caption={Building and running the C++ examples.}, label={lst:tut-run-cpp}]
# configures cpp/examples into build-examples-pixi and builds the all_examples target
pixi run build-examples

# each .cc becomes its own binary; run it FROM THE REPOSITORY ROOT, because the
# YAML-driven examples default to a relative config path such as configs/...
./build-examples-pixi/ex1_propagate_orbit
./build-examples-pixi/ex7_groundstation_odts
./build-examples-pixi/ex7_groundstation_odts configs/my_variant.yaml
\end{lstlisting}

Every YAML-driven example takes an optional first argument overriding the
configuration file, so a scenario variant can be tried without recompiling.

\subsection{Data prerequisites}
\label{sec:tut-running-data}

The Earth-GNSS examples --- \cref{sec:tut-ex3}, \cref{sec:tut-ex5}, and
\cref{sec:tut-ex6} --- are data-driven. They read SP3 precise-ephemeris
products, and in places ANTEX antenna files and RINEX broadcast navigation
messages, which are fetched from NASA CDDIS and require a free Earthdata login
(a \file{~/.netrc} entry, or the \code{EARTHDATA\_USERNAME} and
\code{EARTHDATA\_PASSWORD} environment variables). The download procedure is
specified in \cref{chap:sp3-download}. Downloads are cached under
\file{output/gnss_files/}, so only the first run needs the network.

\begin{lupntwarning}
\code{pixi run build-examples} only \emph{compiles} the programs; the data is
loaded at run time. Each of the three guards its inputs: with the files or
credentials absent it prints an informative note and exits with status zero, so
it always builds and runs, and produces its full analysis once the data is
present. \cref{sec:tut-ex6} additionally requires the staged link and
plasma-delay precompute produced by \file{python/examples/ex6_precompute.py}.
\end{lupntwarning}

\subsection{Where the figures in this chapter come from}
\label{sec:tut-running-figures}

The notebooks are committed \emph{with} their outputs --- \code{nbstripout} is
configured in this repository to preserve them --- so every plot a tutorial
produces already sits in the notebook JSON as a base64 \code{image/png} output.
\file{latex_docs/} is a text-only tree, so those images cannot be committed
beside this chapter. They are extracted instead:

\begin{lstlisting}[language=shell, caption={Extracting the tutorial figures into the data directory.}, label={lst:tut-extract-figs}]
# writes $LUPNT_DATA_PATH/docs/tutorials/<notebook-stem>_<n>.png
LUPNT_DATA_PATH=/path/to/LuPNT_data \
    python latex_docs/scripts/extract_notebook_figures.py

# only one notebook, or a dry run that writes nothing
python latex_docs/scripts/extract_notebook_figures.py --only ex15
python latex_docs/scripts/extract_notebook_figures.py --dry-run
\end{lstlisting}

The counter \code{<n>} runs over that notebook's images in execution order, so
\file{ex1_propagate_orbit_3.png} is the third plot Example~1 draws. The chapter
includes them with \code{\textbackslash datafig}, which degrades to a visible
placeholder when \code{LUPNT\_DATA\_PATH} has not been fetched; the manual
therefore still compiles from a bare clone. Two notebooks contribute no figures:
\cref{sec:tut-ex2} ships without stored outputs (its plots are described in
prose from the plotting code), and \cref{sec:tut-ex13} emits interactive HTML
scenes rather than raster images.

% ===========================================================================
\section{Example 1: Propagating an ELFO Orbit}
\label{sec:tut-ex1}

\textbf{Python:} \file{python/examples/ex1_propagate_orbit.ipynb} \quad
\textbf{C++:} \file{cpp/examples/tutorials/ex1_propagate_orbit.cc}

\paragraph{The scenario.}
The entry point to the library. A lunar Elliptical Frozen Orbit (ELFO) is
initialised at the epoch \textbf{2009-07-15 01:00:00 TDB} from the classical
elements $a=\SI{6541.4}{\kilo\meter}$, $e=0.6$, $i=\SI{56.2}{\degree}$,
$\Omega=\SI{0}{\degree}$, $\omega=\SI{90}{\degree}$, $M_0=\SI{0}{\degree}$,
expressed in the Moon-centred Orbital Plane (OP) frame --- the frame that tracks
the Earth's instantaneous orbit plane about the Moon, and in which the canonical
frozen-orbit parameters of Ely are defined. The Keplerian period of that
semi-major axis is \SI{13.19}{\hour}. Converting to Cartesian and rotating into
the Moon-Centred Inertial (MCI) frame moves the initial position from
$(0,\,1455.58,\,2174.32)\,\si{\kilo\meter}$ in OP to
$(1076.81,\,-334.33,\,2361.16)\,\si{\kilo\meter}$ in MCI --- the same physical
point, a different set of axes.

The force model is a $2\times2$ lunar spherical-harmonic field (so $J_2$ and
$C_{22}$), Earth and Sun as point-mass third bodies, and a cannonball solar
radiation pressure with $C_R=1.8$, $A=\SI{0.02}{\meter\squared}$,
$m=\SI{10}{\kilogram}$. The integrator is adaptive RKF45 with a relative
tolerance of \num{1e-10} and an absolute tolerance of \SI{1e-6}{\meter}. The arc
is \textbf{720 days} --- the variable is spelled \code{sim\_days = 30 * 24} ---
sampled every \SI{30}{\minute}, giving \num{34561} output states.

\begin{lstlisting}[language=py, caption={Elements, frame rotation, and the force model.}, label={lst:tut-ex1-setup}]
# python/examples/ex1_propagate_orbit.ipynb
t0 = pnt.gregorian_to_time(2009, 7, 15, 1, 0, 0)  # [s] TDB

a = 6541.4e3        # [m]   semi-major axis
e = 0.6             # [-]   eccentricity
i = 56.2 * pnt.RAD  # [rad] inclination
raan = 0.0 * pnt.RAD
w = 90.0 * pnt.RAD  # [rad] argument of perigee
M0 = 0.0 * pnt.RAD

coe0_op = np.array([a, e, i, raan, w, M0])
rv0_op = pnt.classical_to_cart(coe0_op, pnt.GM_MOON)                # [m, m/s] in OP
rv0_mci = pnt.convert_frame(t0, rv0_op, pnt.MOON_OP, pnt.MOON_CI)   # [m, m/s] in MCI

dyn = pnt.NBodyDynamics()
dyn.set_frame(pnt.MOON_CI)
dyn.set_units(pnt.SI_UNITS)                    # states in metres / seconds
dyn.add_body(pnt.create_body(pnt.MOON, 2, 2))  # Moon gravity, 2x2 spherical harmonics
dyn.add_body(pnt.create_body(pnt.EARTH))       # Earth third-body
dyn.add_body(pnt.create_body(pnt.SUN))         # Sun third-body + SRP source
dyn.set_srp_coeff(1.8, 0.02, 10.0)             # C_R, area [m^2], mass [kg]
dyn.set_integrator(pnt.IntegratorType.RKF45)
params = pnt.IntegratorParams()
params.reltol = 1e-10
params.abstol = 1e-6                           # [m] -- 1 micrometre position floor
dyn.set_integrator_params(params)
\end{lstlisting}

\paragraph{The computation.}
Propagation is one call. The interpretation happens afterwards: each MCI state
is rotated back into OP and converted to osculating classical elements, i.e.\ the
instantaneous Keplerian orbit that would match the state if all perturbations
were switched off. That turns a Cartesian trajectory into the physically
meaningful $(a,e,i,\omega)$ history in which ``frozen'' means something.

\begin{lstlisting}[language=py, caption={Propagating, then recovering osculating elements in the OP frame.}, label={lst:tut-ex1-prop}]
# python/examples/ex1_propagate_orbit.ipynb
sim_days = 30 * 24
dt_out = 30 * pnt.SECS_MINUTE                                 # [s] output cadence
tspan = np.arange(0, sim_days * pnt.SECS_DAY + dt_out, dt_out)
tfs = t0 + tspan
rv_mci = dyn.propagate(rv0_mci, tfs)                          # [N x 6] in metres, m/s

rv_op = np.array([pnt.convert_frame(tfs[k], rv_mci[k], pnt.MOON_CI, pnt.MOON_OP)
                  for k in range(len(tfs))])
coe = np.array([pnt.cart_to_classical(rv_op[k], pnt.GM_MOON) for k in range(len(tfs))])
# columns: [a, e, i, RAAN, argp, M] -- angles in radians, a in metres
\end{lstlisting}

The last cell re-propagates a short, finely sampled arc (3 orbits at
$T/300 \approx \SI{158}{\second}$ cadence) under an $8\times8$ lunar field and
asks the dynamics object to hand back each acceleration term separately.

\begin{lstlisting}[language=py, caption={Decomposing the acceleration into individual harmonics and third bodies.}, label={lst:tut-ex1-acc}]
# python/examples/ex1_propagate_orbit.ipynb
acc_terms = {}  # term name -> [N, 3] acceleration array [m/s^2]
for k in range(len(tfs_acc)):
    for name, vec in dyn_hf.compute_accelerations(
        tfs_acc[k], rv_acc[k], decompose_gravity=True
    ).items():
        acc_terms.setdefault(name, np.zeros((len(tfs_acc), 3)))[k] = vec

acc_mag = {name: np.linalg.norm(v, axis=1) for name, v in acc_terms.items()}
peak = {name: mag.max() for name, mag in acc_mag.items()}
ranked = sorted(peak, key=peak.get, reverse=True)
\end{lstlisting}

\paragraph{The results.}
The notebook produces four figures.

The first (\cref{fig:tut-ex1-orbit}) is a Plotly 3-D view of the first two
orbital periods ($\approx\SI{26.4}{\hour}$) in the MCI frame, with the Moon
drawn to scale. It is a geometry check: the orbit is visibly eccentric, perilune
sits close to the surface, and the spacecraft spends most of the period out near
apolune, which is exactly the dwell time a relay or navigation satellite is
after.

\begin{figure}[htbp]
  \centering
  \datafig[width=0.52\textwidth]{tutorials/ex1_propagate_orbit_1.png}
  \caption{Example~1, first figure: the ELFO over its first two revolutions in
           the Moon-Centred Inertial frame. The long apolune arc, not the fast
           perilune passage, is where the satellite spends its time.}
  \label{fig:tut-ex1-orbit}
\end{figure}

The second (\cref{fig:tut-ex1-coe}) is the point of the example: a
$2\times2$ panel of the osculating $a$ [\si{\kilo\meter}], $e$ [--],
$i$ [\si{\degree}] and $\omega$ [\si{\degree}] against time in days over the
whole 720-day arc. The semi-major axis oscillates about its mean without a
secular trend; the eccentricity and the argument of perigee execute a slow,
bounded libration about $(e,\omega)=(0.6,\SI{90}{\degree})$ rather than
circulating. That bounded behaviour \emph{is} the frozen-orbit property. (The
panel's suptitle still reads ``30-day'', a leftover from an earlier arc length;
the abscissa runs to 720 days.)

\begin{figure}[htbp]
  \centering
  \datafig[width=0.92\textwidth]{tutorials/ex1_propagate_orbit_2.png}
  \caption{Example~1, second figure: osculating elements in the OP frame over
           the 720-day arc under the N-body plus SRP model. Each element
           oscillates within a bounded envelope; none drifts away.}
  \label{fig:tut-ex1-coe}
\end{figure}

The third (\cref{fig:tut-ex1-phase}) plots the same information as an
$e$-versus-$\omega$ phase portrait, coloured by elapsed time, with a red star at
the nominal frozen point $(\SI{90}{\degree},0.6)$. A truly circulating orbit
would sweep $\omega$ through the full \SI{360}{\degree}; here the trajectory
stays in a closed loop around the star. Reading the colour bar shows the loop is
traversed repeatedly rather than spiralling outward.

\begin{figure}[htbp]
  \centering
  \datafig[width=0.62\textwidth]{tutorials/ex1_propagate_orbit_3.png}
  \caption{Example~1, third figure: the $e$--$\omega$ phase portrait, coloured
           by time. The libration loop around the frozen point (red star) is the
           direct visual signature of a frozen orbit.}
  \label{fig:tut-ex1-phase}
\end{figure}

The fourth (\cref{fig:tut-ex1-acc}) is a semi-logarithmic time history of the
magnitude of each modelled acceleration over three orbits, with dotted vertical
lines at the period boundaries. The ranking it prints establishes the budget
for a lunar navigation satellite: the lunar monopole peaks at
\SI{7.19e-1}{\meter\per\second\squared}; $J_2$ at
\SI{1.56e-4}{\meter\per\second\squared} and the Earth third body at
\SI{1.08e-4}{\meter\per\second\squared} are comparable and four orders of
magnitude down; $C_{31}$ (\SI{5.0e-5}{}) and $C_{22}$ (\SI{2.7e-5}{}) follow;
the Sun contributes \SI{4.6e-7}{\meter\per\second\squared}; and the relativistic
correction (\SI{4.4e-8}{}) and SRP (\SI{3.2e-8}{\meter\per\second\squared}) are
the smallest terms retained. Every harmonic term peaks sharply at perilune and
collapses at apolune, so the higher-order field only matters during the short
low-altitude passage, whereas the Earth's pull peaks at apolune --- the two
dominant perturbations act on opposite halves of the orbit.

\begin{figure}[htbp]
  \centering
  \datafig[width=\textwidth]{tutorials/ex1_propagate_orbit_4.png}
  \caption{Example~1, fourth figure: acceleration budget by force term over
           three orbits with an $8\times8$ lunar field. The lunar harmonics
           spike at perilune; the Earth third body peaks at apolune.}
  \label{fig:tut-ex1-acc}
\end{figure}

The force model, the integrator, and the tolerance contract are specified in
\cref{chap:dynamics} and \cref{chap:integration}; the OP, MCI and body-fixed
frame definitions and the \cpp{ConvertFrame} graph are specified in
\cref{chap:frame-conversions}.

% ===========================================================================
\section{Example 2: Relativistic Time Conversions}
\label{sec:tut-ex2}

\textbf{Python:} \file{python/examples/ex2_time_conversions.ipynb} \quad
\textbf{C++:} \file{cpp/examples/tutorials/ex2_time_conversions.cc}

\paragraph{The scenario.}
The example evaluates the IAU coordinate time scales relevant to cislunar
operations --- $\TT$, $\TDB$, Selenocentric Coordinate Time $\TCL$, and Lunar
Time $\LT$ --- on a \textbf{15-year grid of 721 points from 2020-01-01 to
2035-01-01} (a spacing of about \SI{7.6}{\dayunit}), and a second, denser
\textbf{6-year grid at \SI{0.25}{\dayunit} spacing} used only for the lunar
spectrum. It opens by making the precision argument explicit: the largest
absolute epoch on the grid has an ULP of about \SI{245}{\nano\second}, i.e.\
\SI{73}{\meter} once multiplied by $c$, so the API returns \emph{offsets} rather
than differencing two large absolute epochs.

\paragraph{The computation.}
Three independent routes to $\TT-\TDB$ are compared: the DE440 Eq.~(3)
relativistic integral evaluated from \lupnt{}'s own ephemeris with an \SI{864}{\second}
trapezoid step; JPL's own integrated $\TT-\TDB$ read straight out of the
\code{de440t} time-ephemeris segment (the reference); and the Chebyshev fit of
that kernel, which is what \lupnt{} uses by default.

\begin{lstlisting}[language=py, caption={Three routes to TT - TDB on the same grid.}, label={lst:tut-ex2-routes}]
# python/examples/ex2_time_conversions.ipynb
t_start = float(pnt.gregorian_to_time("2020-01-01T00:00:00"))
t_end = float(pnt.gregorian_to_time("2035-01-01T00:00:00"))
n_pts = 721                                  # ~7.6 day sampling over 15 years
t_tdb = np.linspace(t_start, t_end, n_pts)

# --- Route 2: the DE440t kernel (reference) -----------------------------------
tt_tdb_kernel = np.asarray(
    pnt.spice.get_time_ephemeris_offset(t_tdb, 1000000001), float)   # TT - TDB [s]

# --- Route 1: DE440 Eq. (3) relativistic integral ------------------------------
pnt.set_de440_tt_tdb_step(0.01 * SECS_DAY)   # 864 s trapezoid
tt_tdb_integral = -np.asarray(pnt.tdb_minus_tt_de440(t_tdb), float)

# --- Route 3: Chebyshev fit of the kernel (LuPNT default) ---------------------
pnt.init_tt_minus_tdb_fit(t_start - 30 * SECS_DAY, t_end + 30 * SECS_DAY)
tt_tdb_fit = np.asarray(pnt.tt_minus_tdb(t_tdb), float)
\end{lstlisting}

The expensive $\TDB\leftrightarrow\TCL$ integral --- the one edge in the
conversion graph that sweeps from $T_0=1977$ on every call --- is then composed
with the cheap linear $\TCL\leftrightarrow\LT$ rescaling to reach $\LT-\TT$, and
the composed chain is cross-checked against what \py{Epoch} produces when asked
to walk the graph itself.

\begin{lstlisting}[language=py, caption={Composing LT - TT from offsets, and checking it against Epoch.}, label={lst:tut-ex2-lt}]
# python/examples/ex2_time_conversions.ipynb
tdb_lt = np.asarray(pnt.tdb_minus_lt(t_tdb), float)   # TDB - LT [s]
lt_tt = -tdb_lt - tt_tdb_fit                          # LT - TT = (LT-TDB) + (TDB-TT)

lt_tt_epoch = -np.array([
    float(pnt.time_scale_offset(
        pnt.Epoch.from_seconds(float(x), pnt.Time.LT), pnt.Time.TT))
    for x in (t_tdb[_chk] - tdb_lt[_chk])              # LT reading of the same instant
])

days = (t_tdb - t_tdb[0]) / SECS_DAY
slope_us_per_day = np.polyfit(days, lt_tt * 1e6, 1)[0]
\end{lstlisting}

A final Fourier stage detrends each signal and takes a Hanning-windowed
amplitude spectrum, picking the strongest lines at least three frequency bins
apart. The notebook is explicit that the two signals need different grids: the
\SI{7.6}{\dayunit} spacing of the 15-year record aliases the \SI{27.55}{\dayunit}
anomalistic and \SI{29.53}{\dayunit} synodic months, and the two lunar months
are only resolved once the record is longer than about \SI{4.5}{\yr}, hence the
separate 6-year, \SI{0.25}{\dayunit} grid.

\begin{lstlisting}[language=py, caption={Amplitude spectrum of a detrended, uniformly sampled offset.}, label={lst:tut-ex2-fft}]
# python/examples/ex2_time_conversions.ipynb
def spectrum(sig, t):
    """Amplitude spectrum of a detrended, uniformly-sampled signal."""
    n = len(sig)
    dt = t[1] - t[0]
    x = sig - sig.mean()
    w = np.hanning(n)
    amp = np.abs(np.fft.rfft(x * w)) / (w.sum() / 2.0)
    freq = np.fft.rfftfreq(n, d=dt)
    with np.errstate(divide="ignore"):
        period_d = 1.0 / (freq * SECS_DAY)
    return period_d[1:], amp[1:], freq[1:]
\end{lstlisting}

\paragraph{The results.}
This is the one notebook in the set that is committed \emph{without} stored
outputs, so no extracted figures accompany it; the three figures it draws are
described here from the plotting code.

The first figure is a three-panel stack sharing a year axis. The top panel is
$\TT-\TDB$ from the DE440t kernel in milliseconds, dominated by the annual term
of amplitude $\approx\SI{1.66}{\milli\second}$ driven by the Earth's orbital
eccentricity. The middle panel is the Eq.~(3) integral minus the kernel in
nanoseconds, overlaid with the fitted linear drift: the difference is
\emph{almost purely secular}, so the periodic physics of the two computations
agrees and only the rate differs slightly --- the residual rate comes from the
343 asteroids, 30 KBOs and the Kuiper ring that DE440 integrates and that
\lupnt{} approximates with two ring models. The bottom panel is the Chebyshev fit
minus the kernel in \emph{picoseconds}: the fit interpolates the same data read
as clean offsets, so nothing limits it but the polynomial degree. That is the
argument for making it the default.

The second figure is a two-panel stack for $\LT-\TT$: the top panel shows the
raw offset, whose secular slope is the printed drift compared against the
published \SI{56.0256}{\micro\second\per\dayunit} of Turyshev et al.; the bottom
panel shows what is left after removing that linear trend, a purely periodic
signal at the microsecond level.

The third figure is a pair of log--log amplitude spectra. The left panel, from
the 15-year record, has period in years on the abscissa with a marker at
\SI{1}{\yr}: the annual line towers over everything, with much smaller planetary
lines at longer periods. The right panel, from the dense 6-year record, has
period in days with markers at \SI{27.5545}{\dayunit} (anomalistic month),
\SI{29.5306}{\dayunit} (synodic month) and \SI{13.6061}{\dayunit} (half
anomalistic month); the lunar structure of $\LT-\TT$ resolves into those lines.
The reader's takeaway is the attribution: the annual term of $\TT-\TDB$ is the
Earth's orbit, and the monthly terms of $\LT-\TT$ are the Moon's.

The scale definitions, the conversion graph and its edge classes (constant,
table, linear, model, integral), and the choice of route are specified in
\cref{chap:time-conversions}.

% ===========================================================================
\section{Example 3: GNSS Interface}
\label{sec:tut-ex3}

\textbf{Python:} \file{python/examples/ex3_gnss_interface.ipynb} \quad
\textbf{C++:} \file{cpp/examples/tutorials/ex3_gnss_interface.cc}

\paragraph{The scenario.}
The Earth-GNSS data interfaces, exercised end to end in four parts.

Part~one loads two-line elements archived for \textbf{2025-01-01} ---
\textbf{31 GPS}, \textbf{31 Galileo} and \textbf{4 QZSS} satellites --- and
reports the mean altitude at epoch: \SI{20143}{\kilo\meter} for GPS (range
19\,887--20\,780), \SI{23380}{\kilo\meter} for Galileo (23\,208--25\,728) and
\SI{36622}{\kilo\meter} for QZSS (35\,653--38\,902). Each satellite is then
propagated with two-body dynamics for two days at \SI{1}{\minute} cadence.

Part~two loads transmit-antenna gain tables for four GPS blocks and for Galileo
and QZSS signals. Part~three downloads the daily COD MGEX final SP3 products and
the multi-GNSS broadcast navigation messages for a \textbf{3-day window starting
2026-01-01}, plus one day of lead-in margin, and compares broadcast against
precise orbits and clocks on a \textbf{\SI{15}{\minute} grid of 289 epochs} for
the first five common GPS PRNs (G01--G05) and Galileo PRNs (E02--E06). The SP3
file carries 123 satellites, the broadcast file 115. Part~four evaluates the
satellites' yaw-steering attitude laws.

\begin{lstlisting}[language=py, caption={Downloading and parsing the precise and broadcast products.}, label={lst:tut-ex3-load}]
# python/examples/ex3_gnss_interface.ipynb
START_DT = datetime(2026, 1, 1)   # UTC start date
SIM_DAYS = 3

# one daily product per day, plus a one-day lead-in margin for the interpolators
window_days = [START_DT + timedelta(days=d) for d in range(-1, SIM_DAYS + 1)]
sp3_paths, brdc_paths = [], []
for dt in window_days:
    te = _day_epoch_tai(dt)
    sp3_paths.append(pnt.Sp3Loader.download_file_for_epoch(te, pnt.Time.TAI))
    brdc_paths.append(pnt.RinexNavLoader.download_file_for_epoch(te, pnt.Time.TAI))

sp3 = pnt.Sp3Loader(sp3_paths)          # C++ loaders parse the downloaded files
brdc = pnt.RinexNavLoader(brdc_paths)
antex = pnt.AntexLoader(os.path.join(DATA_DIR, "gnss/igs20.atx"))
\end{lstlisting}

\paragraph{The computation.}
SP3 products are referenced to the satellite centre of mass, while the broadcast
model is used at the antenna phase centre, so the comparison is only meaningful
after applying the IGS20 ANTEX phase-centre offset
$\vec{r}_\mathrm{APC} = \vec{r}_\mathrm{CoM}
 + \mat{R}_{\mathrm{NEU}\to\mathrm{ECEF}}\,\mathbf{PCO}_\mathrm{NEU}$.
The C++ helper builds the satellite body-frame rotation from the nadir direction
and the Sun direction, so the yaw-steering attitude is applied consistently.

\begin{lstlisting}[language=py, caption={Applying the ANTEX phase-centre offset and decomposing into RTN.}, label={lst:tut-ex3-pco}]
# python/examples/ex3_gnss_interface.ipynb
for k in range(len(t_tai_arr)):
    t_k = float(t_tai_arr[k])
    pco_neu = np.asarray(antex.get_pco(const, prn, freq, t_k))   # NEU offset [m]
    r_apc_ecef[k] = np.asarray(
        pnt.AntexLoader.apply_pco_correction_ecef(t_k, r_com_ecef[k], pco_neu))

def rtn_decompose(r_ref, v_ref, dr):
    r_hat = r_ref / np.linalg.norm(r_ref, axis=1, keepdims=True)
    h = np.cross(r_ref, v_ref)
    n_hat = h / np.linalg.norm(h, axis=1, keepdims=True)
    t_hat = np.cross(n_hat, r_hat)
    return (dr * r_hat).sum(axis=1), (dr * t_hat).sum(axis=1), (dr * n_hat).sum(axis=1)
\end{lstlisting}

The clock comparison first removes a systematic per-constellation bias: the
precise and broadcast products are referenced to slightly different time scales,
so their satellite-clock estimates differ by an offset common to every satellite
in a constellation, which a real receiver would absorb into its own clock-bias
state.

\begin{lstlisting}[language=py, caption={Removing the per-constellation broadcast-versus-precise clock bias.}, label={lst:tut-ex3-clkbias}]
# python/examples/ex3_gnss_interface.ipynb
# median(precise - broadcast) over every satellite AND epoch of the constellation
for c in constellations:
    pooled = np.concatenate([
        results[tag]["clk_sp3_raw"] - results[tag]["clk_brdc"]
        for tag in results if const_of(tag) == c])
    clock_bias_s[c] = float(np.nanmedian(pooled))

for tag, d in results.items():
    d["clk_sp3"] = d["clk_sp3_raw"] - clock_bias_s[const_of(tag)]
\end{lstlisting}

\paragraph{The results.}
Ten figures come out of this notebook.

Figure~1 is a 3-D Earth-fixed snapshot of the three constellations propagated
for two days --- GPS in blue, Galileo in orange, QZSS in green --- showing the
two MEO shells and the QZSS inclined-geosynchronous ground trace.

Figures~2 and~3 are the antenna patterns, drawn as polar maps in which the
radius is off-boresight angle from \SI{0}{\degree} at the centre to
\SI{90}{\degree} at the rim, the azimuth is antenna azimuth, and colour is gain
in \si{\dB}. \cref{fig:tut-ex3-gain} shows the four GPS blocks side by side. All
four peak at about \SI{14.8}{\dB} --- IIA, IIR-M and IIF at
$\varphi=\SI{12}{\degree}$, IIR at \SI{10}{\degree} --- and all four show the
main lobe collapsing beyond roughly \SI{14}{\degree}, which is the angular
radius of the Earth as seen from MEO. Everything a lunar receiver can use lives
outside that ring, in the sidelobes, \SI{20}{\dB} to \SI{30}{\dB} down.
Figure~3 puts GPS IIR-M/IIF, Galileo E1/E5a and two QZSS L1 patterns on the same
scale; figure~4 is a 3-D rendering of the Block~IIF table wrapped around
boresight.

\begin{figure}[htbp]
  \centering
  \datafig[width=\textwidth]{tutorials/ex3_gnss_interface_2.png}
  \caption{Example~3, second figure: GPS transmit-antenna gain patterns for the
           four blocks, as polar maps in off-boresight angle (radius, $0$ to
           \SI{90}{\degree}) and antenna azimuth. Peak gain is
           $\approx\SI{14.8}{\dB}$ near \SI{12}{\degree}; a lunar receiver lives
           entirely in the sidelobe structure outside the main beam.}
  \label{fig:tut-ex3-gain}
\end{figure}

Figure~5 (\cref{fig:tut-ex3-rtn}) is the orbit comparison: six panels of
broadcast minus (precise + PCO) in the radial, along-track and cross-track
directions, GPS in the left column and Galileo in the right, over the 72-hour
window. The per-PRN standard deviations printed alongside are, for GPS,
$\sigma_R = 0.14$--\SI{0.19}{\meter}, $\sigma_T = 0.48$--\SI{0.65}{\meter} and
$\sigma_N = 0.21$--\SI{0.42}{\meter}; for Galileo they are roughly half that,
$\sigma_R = 0.12$--\SI{0.16}{\meter}, $\sigma_T = 0.26$--\SI{0.35}{\meter},
$\sigma_N = 0.11$--\SI{0.20}{\meter}. The along-track channel is the worst in
both systems, and the traces show the characteristic once-per-revolution
signature of a broadcast fit.

\begin{figure}[htbp]
  \centering
  \datafig[width=\textwidth]{tutorials/ex3_gnss_interface_5.png}
  \caption{Example~3, fifth figure: broadcast minus (precise + phase-centre
           offset) orbit differences in the radial/along-track/cross-track
           frame, GPS (left) and Galileo (right), \SI{15}{\minute} cadence over
           three days. Along-track is the largest component in both systems.}
  \label{fig:tut-ex3-rtn}
\end{figure}

Figure~6 (\cref{fig:tut-ex3-clk}) is the clock comparison, again in a
$3\times2$ grid: the broadcast clock polynomial, the debiased precise clock, and
their difference in nanoseconds. The removed per-constellation bias is
\SI{0.789}{\nano\second} for GPS and \SI{-1.005}{\nano\second} for Galileo. What
remains is broadcast clock \emph{prediction} error: per-PRN standard deviations
of 0.32 to \SI{1.63}{\nano\second} with peak-to-peak excursions up to
\SI{8.4}{\nano\second}, which at $c$ is a few metres of pseudorange.

\begin{figure}[htbp]
  \centering
  \datafig[width=0.86\textwidth]{tutorials/ex3_gnss_interface_6.png}
  \caption{Example~3, sixth figure: broadcast clock, debiased precise clock, and
           their difference for five GPS and five Galileo satellites, after
           removing the per-constellation median offset.}
  \label{fig:tut-ex3-clk}
\end{figure}

Figures~7 and~8 pool those differences into histograms, one row per
constellation, with radial, along-track, cross-track and clock panels annotated
with RMS and $\sigma$. Pooled over 5 PRNs $\times$ 289 epochs the GPS numbers are
radial RMS \SI{0.53}{\meter} (95th percentile \SI{0.96}{\meter}), along-track
\SI{0.82}{\meter} (\SI{1.62}{\meter}), cross-track \SI{0.32}{\meter}
(\SI{0.60}{\meter}) and clock \SI{1.02}{\nano\second} (\SI{1.75}{\nano\second});
the Galileo numbers are \SI{0.16}{\meter}, \SI{0.31}{\meter}, \SI{0.16}{\meter}
and \SI{0.84}{\nano\second}. This is the signal-in-space error budget that the
later measurement simulations must tolerate --- and it is dominated by the
broadcast product, not by anything \lupnt{} models.

Figure~9 compares nominal and block-specific ``dedicated'' yaw-steering laws for
one GPS and one Galileo satellite over the window. The two curves \emph{coincide}
and their difference is flat, which is the expected result and is explained
rather than hidden: over this arc G01 has $|\beta| =$ 23.8--\SI{25.1}{\degree},
far outside the \SI{5.8}{\degree} eclipse regime in which a rate-limited law
departs from nominal. Figure~10 (\cref{fig:tut-ex3-yaw}) therefore constructs a
synthetic low-$\beta$ turn at $\beta=\SI{1.5}{\degree}$ and sweeps the orbit
angle $\pm\SI{40}{\degree}$ around noon and around midnight, showing the nominal
law against the constant-rate block laws (IIR and IIR-M at
\SI{0.20}{\degree\per\second}, IIF at \SI{0.11}{\degree\per\second}). Only there
do the laws visibly diverge, and the slower IIF limit lags furthest behind
nominal.

\begin{figure}[htbp]
  \centering
  \datafig[width=\textwidth]{tutorials/ex3_gnss_interface_10.png}
  \caption{Example~3, tenth figure: a synthetic low-$\beta$ noon turn (left) and
           midnight turn (right) at $\beta=\SI{1.5}{\degree}$, comparing the
           nominal yaw law with the rate-limited GPS block laws.}
  \label{fig:tut-ex3-yaw}
\end{figure}

The observable definitions, the antenna gain model, and the broadcast ephemeris
evaluation this example compares against are specified in
\cref{chap:gnss-measurements}; the download workflow and product naming are in
\cref{chap:sp3-download}.

% ===========================================================================
\section{Example 4: Plasmasphere and Ionosphere}
\label{sec:tut-ex4}

\textbf{Python:} \file{python/examples/ex4_plasmasphere.ipynb} \quad
\textbf{C++:} \file{cpp/examples/tutorials/ex4_plasmasphere.cc}

\paragraph{The scenario.}
The \code{pecsim} backend behind \py{pylupnt.plasma}, in two parts, both at the
epoch \textbf{2026-01-01 00:00 UTC} ($t=\SI{820497627}{\second}$ from J2000)
with a geomagnetic activity index $K_p = 2.0$ and the IRI-2007 ionosphere model
selected inside GCPM v2.4.

Part~one samples the electron density on a $101\times101$ meridional grid
spanning $\pm 8\,R_E$ in the noon--midnight plane, skipping points inside
$1.01\,R_E$ and outside the domain. Part~two ray-traces one GPS-to-lunar link:
the transmitter is a circular MEO satellite at
$a=\SI{26560}{\kilo\meter}$, $e=0.001$, $i=\SI{55}{\degree}$ with its mean
anomaly deliberately tuned to \SI{245}{\degree} so that it sits on the
\emph{far} side of the Earth from the Moon; the receiver is the same ELFO family
used throughout ($a=\SI{6541.4}{\kilo\meter}$, $e=0.6$, $i=\SI{56.2}{\degree}$,
$\omega=\SI{90}{\degree}$, $M=\SI{180}{\degree}$), propagated about the Moon and
then shifted into the Earth-centred frame. At this epoch the transmitter is at
$4.17\,R_E$, the Moon at $56.67\,R_E$ and the receiver at $55.34\,R_E$, and the
line of sight grazes the plasmasphere instead of missing it entirely. The
notebook says so explicitly: a transmitter on the same side as the Moon would
give a path that never dips below the $4\,R_E$ cutoff and a TEC of exactly zero.

\begin{lstlisting}[language=py, caption={Epoch, activity index, and the electron-density grid.}, label={lst:tut-ex4-grid}]
# python/examples/ex4_plasmasphere.ipynb
YEAR, MONTH, DAY, HOUR = 2026, 1, 1, 0
mjd = plasma.gregorian_to_mjd(YEAR, MONTH, DAY, HOUR)
epoch_j2000 = plasma.mjd_to_tj2000(mjd)   # seconds since J2000 (UTC)
dt = plasma.mjd_to_datetime(mjd)          # DateTime for GCPM/IRI queries
kp = plasma.get_kp_index(dt)
plasma.set_iri_model("IRI2007")

for i in range(N):
    for j in range(N):
        x, z = XX[i, j], ZZ[i, j]
        r = np.hypot(x, z)
        alatr = np.arctan2(z, np.abs(x))     # magnetic latitude [rad]
        amlt = 12.0 if x >= 0 else 0.0       # noon-midnight meridian
        ne = plasma.gcpm_v24(dt, r, amlt, alatr, kp)[0]
\end{lstlisting}

\paragraph{The computation.}
The ray tracer integrates the plasma model along the path at
\SI{100}{\kilo\meter} steps with RK4, applies a receiver-position correction by
Nelder--Mead, and stops modelling plasma beyond a cutoff radius of $4\,R_E =
\SI{25484}{\kilo\meter}$ (vacuum above that, all the way to the Moon).

\begin{lstlisting}[language=py, caption={Ray-trace configuration and the returned path profile.}, label={lst:tut-ex4-trace}]
# python/examples/ex4_plasmasphere.ipynb
config = plasma.RayTraceConfig()
config.freq_Hz = plasma.freq_L1        # 1575.42 MHz
config.step_size = 100.0               # integration step [km]
config.correction = True               # apply receiver-position correction
config.fine_correction = False
config.correction_method = "neldermead"
config.cutoff_r = CUTOFF_RE * RE_km    # 4 R_E
config.kp = kp
config.integ_method = "rk4"
config.straight_ray = False

pp = plasma.trace_ray(epoch_j2000, pos_tx_km, pos_rx_km, config,
                      debug_prop=False, debug_corr=False)
# pp.tecu, pp.tec_delay_m, pp.dist_bend_m, pp.dist_straight_km, pp.pos_eci, pp.r, pp.s
\end{lstlisting}

\paragraph{The results.}
The example prints a small table of scalars and draws two figures.

The scalars are the physical answer. The grid sweep reports an electron density
spanning \SI{1.58e-1}{} to \SI{1.55e6}{\per\centi\meter\cubed} --- seven orders
of magnitude between the outer magnetosphere and the F-region peak. The ray
trace returns \textbf{\SI{4.968}{\TECU}} of total electron content along a
straight-line distance of \SI{376900}{\kilo\meter}, which at L1 becomes a total
delay of \SI{0.8066}{\meter}, of which \SI{0.8067}{\meter} is the dispersive TEC
term and only \SI{2e-6}{\meter} is geometric bending. The lesson is quantitative:
the delay is dominated entirely by the TEC integral, ray bending is negligible at
L1, and the whole effect is at the metre level --- large compared with a
carrier-phase observable, small compared with a broadcast ephemeris error.

The first figure (\cref{fig:tut-ex4-slice}) is the meridional slice: a
\code{pcolormesh} of $\log_{10} n_e$ over $X\in[-8,8]\,R_E$ (noon on $+X$,
midnight on $-X$) and $Z\in[-8,8]\,R_E$ along the magnetic axis, colour-scaled
from $10^0$ to $10^5\,\si{\per\centi\meter\cubed}$, with the Earth filled in, a
dashed white plasmapause circle at $4\,R_E$ and the GPS orbit radius
($4.2\,R_E$) drawn in yellow. The dense torus hugging the Earth is the
plasmasphere; the sharp drop across the plasmapause and the day/night asymmetry
between the noon and midnight halves are both visible, and the GPS shell is
shown to sit right at the plasmapause --- which is why link geometry, not
altitude alone, decides whether a signal is affected.

\begin{figure}[htbp]
  \centering
  \datafig[width=0.68\textwidth]{tutorials/ex4_plasmasphere_1.png}
  \caption{Example~4, first figure: meridional slice of GCPM~v2.4 + IRI-2007
           electron density at $K_p=2$, in the noon--midnight plane. Dashed
           white: the $4\,R_E$ plasmapause; yellow: the GPS orbit radius.}
  \label{fig:tut-ex4-slice}
\end{figure}

The second figure (\cref{fig:tut-ex4-ray}) has three panels. The left panel
plots the traced ray in the $XZ$ plane over the near-Earth region, each sample
coloured by $\log_{10} n_e$, with the transmitter marked, the $4\,R_E$ cutoff
circle dashed, and a cyan arrow pointing off-plot toward the Moon; it shows the
path diving in toward perigee and back out. The middle panel is altitude above
the Earth's surface against arc length, truncated at the outbound $4\,R_E$
crossing, with the perigee marked. The right panel is the electron-density
profile along the same arc. The notebook then spends a whole markdown cell on
that third panel: the two near-vertical walls in $n_e$ are \emph{not} a
discontinuity but the genuinely steep, still-continuous plasmapause gradient,
under-sampled by the \SI{100}{\kilo\meter} step; resampling the identical path
at \SI{2}{\kilo\meter} shows $n_e$ falling smoothly from about 2000 to
\SI{330}{\per\centi\meter\cubed} over \SI{230}{\kilo\meter} on the inbound wall
and rising from about 450 to \SI{4400}{\per\centi\meter\cubed} over
\SI{450}{\kilo\meter} on the perigee side. The density peak also sits slightly
off the geometric perigee, because the plasma is organised by dipole $L$-shell
and magnetic latitude rather than by geocentric radius.

\begin{figure}[htbp]
  \centering
  \datafig[width=\textwidth]{tutorials/ex4_plasmasphere_2.png}
  \caption{Example~4, second figure: the GPS-to-ELFO ray through the
           plasmasphere (left, coloured by $n_e$), its altitude profile to the
           $4\,R_E$ crossing (centre), and the electron density along the path
           (right).}
  \label{fig:tut-ex4-ray}
\end{figure}

The GCPM and IRI density models, the Haselgrove ray-tracing equations, the
dispersive delay relation, and the model boundaries are specified in
\cref{chap:ionosphere}. A higher-fidelity closed-loop variant with a full
constellation and occultation gating lives in
\file{projects/old/Plasma_Examples/ex_plasma_raytrace.py}.

% ===========================================================================
\section{Example 5: GNSS Measurement Simulation}
\label{sec:tut-ex5}

\textbf{Python:} \file{python/examples/ex5_gnss_measurement_sim.ipynb} \quad
\textbf{C++:} \file{cpp/examples/tutorials/ex5_gnss_measurement_sim.cc}

\paragraph{The scenario.}
Which Earth-GNSS signals can a lunar receiver actually track? The receiver is
the standard ELFO ($a=\SI{6541.4}{\kilo\meter}$, $e=0.6$,
$i=\SI{56.2}{\degree}$, $\omega=\SI{90}{\degree}$), started at
\textbf{2026-01-01 00:00:00 TDB} and propagated for exactly \textbf{one orbital
period (\SI{47474.9}{\second}, \SI{13.19}{\hour})} at \SI{300}{\second} steps,
i.e.\ 159 receiver epochs. Over that period the altitude above the lunar surface
swings from \SI{879}{\kilo\meter} at periapsis to \SI{8727}{\kilo\meter} at
apoapsis.

The transmitters come from the same COD MGEX SP3 products as
\cref{sec:tut-ex3}: \textbf{32 GPS} and \textbf{27 Galileo} satellites with
complete coverage over the arc, sampled on a \SI{300}{\second} grid that starts
one step early so the light-time iteration (about \SI{1.3}{\second} at cislunar
range) has data to interpolate. The receiver antenna is the \code{moongpsr}
model --- an Earth-pointing dish, \SI{+14}{\dBi}, roughly \SI{7}{\degree}
half-power beamwidth --- and both the Earth and the Moon are registered as
occulting bodies, the Moon with a per-epoch position provider because it moves
\SI{7}{\degree} over the arc. Acquisition and tracking $C/N_0$ thresholds are
\SI{22}{\dBHz} and \SI{20}{\dBHz}.

\begin{lstlisting}[language=py, caption={Receiver-side measurement model configuration.}, label={lst:tut-ex5-opts}]
# python/examples/ex5_gnss_measurement_sim.ipynb
opts = pnt.GnssMeasurementOptions()
opts.frame = pnt.ECI
opts.receive_time_scale = pnt.TAI
opts.solve_light_time = True
opts.apply_transmitter_relativity = True
opts.apply_visibility = True
opts.apply_cn0_threshold = True
opts.cn0_acquisition_threshold_dbhz = 22.0
opts.cn0_tracking_threshold_dbhz = 20.0

gnss_meas = pnt.GNSSMeasurements(gps_const)
gnss_meas.set_frequency(pnt.GnssFreq.L1)
gnss_meas.add_constellation(gal_const, pnt.GnssFreq.E1)
gnss_meas.set_options(opts)
gnss_meas.set_receiver_antenna(pnt.Antenna("moongpsr"))  # Earth-pointing, +14 dBi
gnss_meas.set_occluding_bodies([earth_body, moon_body])
gnss_meas.set_boresight_target_provider(earth_pos_eci)
\end{lstlisting}

\paragraph{The computation.}
For each receiver epoch the model solves light time, applies transmitter
relativity, checks the path against both occulting bodies, evaluates the
transmit antenna gain in the satellite body frame, adds free-space loss, and
returns one channel per surviving link with its $C/N_0$.

\begin{lstlisting}[language=py, caption={Building the per-epoch channel list.}, label={lst:tut-ex5-loop}]
# python/examples/ex5_gnss_measurement_sim.ipynb
epoch_results = []
for t_tai, state in zip(t_tai_arr.tolist(), states_list):
    channels = gnss_meas.build_channels(t_tai, state)
    epoch_results.append(_Epoch(t_tai, channels))

# Use ch.gnss_const (not the PRN number) to tag constellation -- Galileo PRNs
# 2-31 overlap with GPS PRNs, so PRN-based tagging silently miscounts.
for i, ep in enumerate(epoch_results):
    for ch in ep.channels:
        const = "G" if ch.gnss_const == pnt.GnssConst.GPS else "E"
        links.append((t_hrs[i], f"{const}{ch.prn:02d}", float(ch.cn0_dbhz)))
\end{lstlisting}

Each visible link is then projected back into the \emph{transmitter's} antenna
frame, so the reader can see where on the pattern the lunar receiver actually
sits.

\begin{lstlisting}[language=py, caption={Projecting each link into the transmit-antenna frame.}, label={lst:tut-ex5-angles}]
# python/examples/ex5_gnss_measurement_sim.ipynb
att = pnt.GnssAttitude(tx_pos, sun_pos)
u = rx_pos - tx_pos
u /= np.linalg.norm(u)
theta_r, phi_r = att.get_angles(np.array(u))   # theta = azimuth, phi = off-boresight
\end{lstlisting}

\paragraph{The results.}
Over the 159 epochs the model yields \textbf{844 visible links} ---
\textbf{657 GPS} and \textbf{187 Galileo} --- a mean of \textbf{5.3 visible
satellites}, a peak of \textbf{10} and a minimum of \textbf{1}, with $C/N_0$
spanning \SI{20.0}{} to \SI{44.9}{\dBHz}. Four or more simultaneous satellites,
the classic requirement for an instantaneous 3-D-plus-clock fix, is available
much but not all of the time --- which is precisely why \cref{sec:tut-ex6} uses
a filter rather than a snapshot solution.

The first figure (\cref{fig:tut-ex5-vis}) is a two-panel step plot of the number
of visible GPS (top) and Galileo (bottom) satellites against time from epoch in
hours. The dips are geometry: occultation by the Moon near perilune, and stretches
where the Earth-pointing beam sees only weak sidelobes. Galileo contributes
roughly a quarter of the links, reflecting both the smaller usable sidelobe
structure and the smaller number of satellites in the SP3 file.

\begin{figure}[htbp]
  \centering
  \datafig[width=0.92\textwidth]{tutorials/ex5_gnss_measurement_sim_1.png}
  \caption{Example~5, first figure: visible GPS (top) and Galileo (bottom)
           satellite count over one ELFO period, gated at \SI{22}{\dBHz}
           acquisition.}
  \label{fig:tut-ex5-vis}
\end{figure}

The second figure (\cref{fig:tut-ex5-cn0}) is the more informative view: a
$C/N_0$ heat map with satellites on the ordinate (GPS then Galileo, PRN-sorted)
and time on the abscissa, colour-scaled \SI{15}{} to \SI{45}{\dBHz}. Each
horizontal streak is one satellite's pass; the streaks are short, because a lunar
receiver sees a given GPS satellite only while the sidelobe geometry happens to
line up. The colour tells the reader what the visibility count cannot: a link
barely over the \SI{22}{\dBHz} threshold and a \SI{40}{\dBHz} link contribute
very different measurement noise.

\begin{figure}[htbp]
  \centering
  \datafig[width=\textwidth]{tutorials/ex5_gnss_measurement_sim_2.png}
  \caption{Example~5, second figure: $C/N_0$ received by the ELFO from each GPS
           L1 and Galileo E1 satellite over one orbit. Short streaks are
           sidelobe passes, not continuous tracking.}
  \label{fig:tut-ex5-cn0}
\end{figure}

Figures~3 and~4 close the argument. Each is a $2\times3$ grid of polar plots,
one per most-tracked PRN (GPS 2, 11, 15, 18, 20, 29; Galileo 14, 18, 21, 23, 27,
30), showing that transmitter's own gain table as a faint background with every
link to the lunar receiver over-plotted at its off-boresight angle and azimuth.
\cref{fig:tut-ex5-polar} shows the GPS panel. Essentially every point lies well
outside the Earth's angular radius from MEO ($\approx\SI{13.9}{\degree}$), i.e.\
outside the main beam that the antenna was designed for. That is the whole thesis
of lunar GNSS in one picture: the usable signal is sidelobe signal.

\begin{figure}[htbp]
  \centering
  \datafig[width=\textwidth]{tutorials/ex5_gnss_measurement_sim_3.png}
  \caption{Example~5, third figure: directions to the lunar receiver plotted on
           each GPS transmitter's own gain pattern. Every link falls outside the
           Earth-facing main lobe.}
  \label{fig:tut-ex5-polar}
\end{figure}

The light-time iteration, the relativistic and plasma corrections, the antenna
gain and occultation gating, and the online versus batch \cpp{Precompute} paths
are specified in \cref{chap:gnss-measurements}; the $C/N_0$ computation itself
is in \cref{chap:link-budget}.

% ===========================================================================
\section{Example 6: GNSS Sidelobe ODTS}
\label{sec:tut-ex6}

\textbf{Python:} \file{python/examples/ex6_gnss_odts.ipynb} \quad
\textbf{C++:} \file{cpp/examples/tutorials/ex6_gnss_odts.cc}

\paragraph{The scenario.}
Orbit determination and time synchronisation for a lunar receiver navigating on
the weak sidelobe signals of \cref{sec:tut-ex5}. The scenario is a YAML file,
\file{configs/lunar_gnss_odts.yaml}: a physical \code{receiver} \cpp{Spacecraft}
agent hosts a \cpp{LunarGnssOdtsApp}, and the notebook's own configuration
module mirrors the same values field-for-field so that the C++ precompute cache
--- keyed on a fingerprint of the configuration --- is shared between the
notebook and the standalone script. The configuration the notebook prints is:

\begin{itemize}[nosep]
  \item a \textbf{\SI{13.19}{\hour}} run ($\approx$ one ELFO period,
        \SI{47475}{\second}) at \textbf{\SI{1}{\hertz}}, starting
        \textbf{2026-01-01T02:00:00};
  \item TDCP on, SRP-coefficient estimation on with a truth
        $C_R A/m = \SI{2.00e-3}{\meter\squared\per\kilogram}$;
  \item ionosphere-free (L1+L5) pseudorange, $C/N_0$-derived noise plus filter
        inflation of \SI{10}{\meter} on pseudorange and \SI{200}{\milli\meter}
        on TDCP;
  \item tangent-altitude cutoffs of \SI{1000}{\meter} for pseudorange and
        \SI{4000}{\meter} for TDCP;
  \item GPS + Galileo, \code{moongpsr} receiver antenna, \SI{22}{}/\SI{20}{\dBHz}
        acquisition/tracking gate;
  \item plasma delay ray-traced every \SI{120}{\second} and interpolated to
        \SI{1}{\hertz}.
\end{itemize}

The filter is deliberately imperfect: truth carries higher-fidelity dynamics,
the Sun's Shapiro delay and the ray-traced plasmaspheric delay, while the
estimator uses a simpler onboard-style model and inflated measurement noise so
that it does not chase unmodelled effects as if they were state error.

\paragraph{The computation.}
Three observables drive a stochastic-cloning UDU extended Kalman filter over the
nine-state $[\vec{r},\vec{v},b,\dot b, C_R A/m]$: pseudorange, instantaneous
Doppler, and time-differenced carrier phase. TDCP references two epochs, so the
augmented state carries both the current and the previous receiver state --- the
main mathematical lesson of the notebook.

Stage~1 builds the link geometry and caches it; stage~2 ray-traces the plasma
delay. Ray tracing is far too slow to run per link at \SI{1}{\hertz}, so it is
done at \SI{120}{\second} anchor epochs and interpolated in time per
satellite-frequency link, cutting the trace count by roughly $120\times$; a
satellite whose whole pass falls between two anchors is traced directly rather
than silently zeroed.

\begin{lstlisting}[language=py, caption={Loading the scenario, then reading back the two cached precompute stages.}, label={lst:tut-ex6-stage}]
# python/examples/ex6_gnss_odts.ipynb
cfg = build_config()     # struct config: drives the Stage 2 plasma ray-trace
scen = load_scenario()   # agent-based scenario dict for pnt.Simulation
sim = pnt.Simulation(scen)
app = sim.get_agent("receiver").get_application()

app.precompute()         # Stage 1: link geometry / C-N0 -> precomputed_links.csv
links = pd.read_csv(cfg.links_file)

if delays_cache_valid(cfg, PLASMA_DELAY_DT_S):     # Stage 2: GCPM/IRI ray trace
    delays = pd.read_csv(Path(cfg.delays_file))
else:
    delays = compute_delays(cfg, links, PLASMA_DELAY_DT_S)
\end{lstlisting}

\begin{lstlisting}[language=py, caption={Running the filter and reading the per-run summary.}, label={lst:tut-ex6-run}]
# python/examples/ex6_gnss_odts.ipynb
sim.run()
summary = app.get_summaries()[0]
print(f"Final position error:    {summary.final_position_error_m:8.3f} m")
print(f"RMS position error:      {summary.rms_position_error_m:8.3f} m")
print(f"Final clock bias error:  {summary.final_clock_bias_error_m * 1e2:8.3f} cm")

traj = pd.read_csv(Path(cfg.output_dir) / "trajectory_mc0.csv")
\end{lstlisting}

The Monte-Carlo section re-runs the same deterministic truth and link geometry
with only the noise realisation and initial-estimate perturbation changing,
three trials each for pseudorange-only and pseudorange-plus-TDCP:

\begin{lstlisting}[language=py, caption={A three-trial Monte Carlo for each measurement variant.}, label={lst:tut-ex6-mc}]
# python/examples/ex6_gnss_odts.ipynb
N_MC = 3
MC_VARIANTS = {"Pseudorange only": ("mc_pr_only", False),
               "Pseudorange + TDCP": ("mc_pr_tdcp", True)}

s = load_scenario()
a = s["agents"]["receiver"]["application"]
a["simulation"]["output_dir"] = str(vdir)
a["simulation"]["monte_carlo_runs"] = N_MC
a["measurements"]["use_tdcp"] = use_tdcp
pnt.Simulation(s).run()
\end{lstlisting}

\paragraph{The results.}
Ten figures, and they build an argument rather than just reporting numbers.

Figure~1 (\cref{fig:tut-ex6-delay}) is the plasma sanity check: a scatter of
plasmaspheric delay (logarithmic, in metres) against the ray's tangent altitude
above the Earth's surface in kilometres, coloured by epoch index, for the L1
links only. The delay falls off monotonically as the tangent point rises away
from the dense inner plasmasphere, which is the expected physical behaviour; the
delays span \SI{0}{} to \SI{133.91}{\meter} across the arc.

\begin{figure}[htbp]
  \centering
  \datafig[width=0.78\textwidth]{tutorials/ex6_gnss_odts_1.png}
  \caption{Example~6, first figure: GCPM/IRI-2007 plasmaspheric delay against
           signal-path tangent altitude, L1 links. Delay falls off
           monotonically as the path clears the plasmasphere.}
  \label{fig:tut-ex6-delay}
\end{figure}

Figure~2 (\cref{fig:tut-ex6-rtn}) is the filter's core result: six panels of
estimate error against time in hours, radial/tangential/normal position in the
left column (clipped to $\pm\SI{100}{\meter}$) and the corresponding velocity
components in the right column (clipped to $\pm\SI{20}{\milli\meter\per\second}$),
each overlaid with the filter's own formal $3\sigma$ envelope from the EKF
covariance. The error stays inside the envelope, and the envelope narrows as
sidelobe measurements accumulate.

\begin{figure}[htbp]
  \centering
  \datafig[width=\textwidth]{tutorials/ex6_gnss_odts_2.png}
  \caption{Example~6, second figure: RTN position and velocity error with the
           filter's formal $3\sigma$ bounds, over one ELFO period at
           \SI{1}{\hertz}.}
  \label{fig:tut-ex6-rtn}
\end{figure}

Figure~3 shows the clock: bias error and drift error in range-equivalent units
with their $3\sigma$ bands, the drift band explicitly flagged as untrustworthy
(a covariance-conditioning artefact, unrelated to TDCP). Figure~4 is the
solar-radiation-pressure coefficient estimate relative to truth, with its
$3\sigma$ band: the filter starts from a deliberately biased value and the
estimate creeps toward the dashed truth line while the band shrinks slowly ---
SRP acts as a slow, mostly along-track acceleration, so one orbital period is
about the shortest arc on which it moves at all.

Figure~5 (\cref{fig:tut-ex6-hist}) turns the time histories into four
histograms over all 47\,475 epochs of the single run. The printed statistics are
radial mean \SI{-3.43}{\meter}, $\sigma=\SI{20.48}{\meter}$, RMS
\SI{20.76}{\meter}, 95th-percentile $|{\cdot}|$ \SI{32.84}{\meter}; transverse
\SI{5.24}{}/\SI{24.25}{}/\SI{24.81}{}/\SI{47.33}{\meter}; normal
\SI{-5.23}{}/\SI{22.28}{}/\SI{22.88}{}/\SI{51.63}{\meter}; and clock bias
\SI{-6.60}{}/\SI{20.32}{}/\SI{21.37}{}/\SI{33.95}{\meter}. In other words, a
lunar receiver navigating purely on terrestrial GNSS sidelobes holds a few tens
of metres in each axis over an orbit.

\begin{figure}[htbp]
  \centering
  \datafig[width=\textwidth]{tutorials/ex6_gnss_odts_5.png}
  \caption{Example~6, fifth figure: per-epoch distribution of the RTN position
           error and the range-equivalent clock-bias error. Each is centred near
           zero with an RMS in the low tens of metres.}
  \label{fig:tut-ex6-hist}
\end{figure}

Figure~6 overlays, on twin axes, the number of distinct GNSS satellites tracked
per epoch and the total number of measurement rows fed to the filter (pseudorange
and Doppler rows plus one TDCP row per continuously tracked link); 47\,474 of
47\,475 epochs carry at least one TDCP row, the exception being the first, which
has no previous epoch to difference against. Figure~7 histograms the broadcast
signal-in-space error the pseudorange actually sees, split into the transmitter
orbit error projected on the line of sight, the transmitter clock error in range
units, and their sum. Pooled over 78\,391 GPS satellite-epochs the total is
mean \SI{+0.447}{\meter}, RMS \SI{0.726}{\meter}, 95th percentile
\SI{1.197}{\meter}; over 30\,315 Galileo satellite-epochs it is
\SI{+0.176}{\meter}, RMS \SI{0.319}{\meter}, 95th percentile \SI{0.628}{\meter}
--- consistent with \cref{sec:tut-ex3} and, the notebook argues, the binding
residual in this scenario.

Figures~8 and~9 are the Monte-Carlo comparison. Figure~8 plots the 3-D position
error of each of the three trials plus the ensemble mean, pseudorange-only on
the left and pseudorange-plus-TDCP on the right, on a common
\SIrange{0}{200}{\meter} scale. \cref{fig:tut-ex6-mc} is the head-to-head bar
chart: mean RMS position error drops from \SI{47.27}{\meter} to
\SI{40.10}{\meter} and mean final error from \SI{21.24}{\meter} to
\SI{19.16}{\meter} when TDCP is added, with the individual trials shown as black
dots so the reader can see the spread is comparable to the improvement.

\begin{figure}[htbp]
  \centering
  \datafig[width=0.82\textwidth]{tutorials/ex6_gnss_odts_9.png}
  \caption{Example~6, ninth figure: RMS and final 3-D position error,
           pseudorange-only versus pseudorange-plus-TDCP, three Monte-Carlo
           trials each (bars are ensemble means, dots individual trials).}
  \label{fig:tut-ex6-mc}
\end{figure}

The accompanying discussion is the sharpest quantitative result in the notebook
and does not appear in any figure. TDCP is single-frequency (L1 / E1), so unlike
the ionosphere-free pseudorange it gets \emph{no} plasma cancellation; the
epoch-to-epoch change of the ray-traced delay, measured directly from the cached
table, is RMS \SI{3.6}{\centi\meter} with a tail to about \SI{0.85}{\meter}.
Feeding TDCP at a nominal \SI{1}{\centi\meter} noise makes the filter chase those
spikes: normalised-error RMS of 25/31/49 in R/T/N and a position RMS of
\SI{768}{\meter}, i.e.\ divergence. At the default \SI{0.20}{\meter} inflation
the normalised-error RMS is 1.6 in all three axes --- as consistent as the
pseudorange-only baseline --- \emph{and} more accurate. Measurement-noise
inflation, not process noise, is the lever.

Figure~10 is a capability demonstration rather than an accuracy result: the same
engine run at the future epoch \textbf{2027-01-01T02:00:00} for
\SI{15}{\minute}, GPS-only, with the constellation seeded from a freshly
downloaded YUMA almanac and numerically propagated (Earth $J_2$ plus Sun and
Moon third bodies) to the run epoch, and with a modelled synthetic
signal-in-space error injected through the same filter path. Over 901 epochs it
gives a final 3-D position error of \SI{8522}{\meter} and an RMS of
\SI{18665}{\meter} --- kilometre-level, as the notebook says it should be for a
year-extrapolated almanac seen through sparse sidelobes. The point is that the
pipeline runs and produces sane measurements at an epoch for which no precise
product exists.

Run \file{python/examples/ex6_precompute.py} first: stage~1 builds the link
geometry, stage~2 the GCPM/IRI plasma delays, and both caches are keyed on a
fingerprint of the configuration so the notebook and the C++ binary share them.
The filter formulation is specified in \cref{chap:filters}, the observables in
\cref{chap:gnss-measurements}, and the delay model in \cref{chap:ionosphere}.

% ===========================================================================
\section{Example 7: Ground-Station Orbit Determination}
\label{sec:tut-ex7}

\textbf{Python:} \file{python/examples/ex7_groundstation_odts.ipynb} \quad
\textbf{C++:} \file{cpp/examples/tutorials/ex7_groundstation_odts.cc}

\paragraph{The scenario.}
A lunar satellite is tracked by the three \SI{70}{\meter} Deep Space Network
antennas --- Goldstone (\code{DSS14}), Canberra (\code{DSS43}) and Madrid
(\code{DSS63}) --- via two-way range and range rate over a
\textbf{\SI{36}{\hour}} arc sampled at \textbf{433 epochs}, producing
\textbf{627 measurements} in total under a \SI{10}{\degree} elevation mask.
This is the first fully agent-based example: a shared \code{World} holds the
epoch, frame and truth force model; a \cpp{Satellite} agent self-propagates the
ELFO truth; each antenna is a \cpp{GroundStation} agent running a
\cpp{GroundStationTrackingApp} that generates its own visibility-gated, noisy
observations; and a \cpp{GroundStationManager} agent runs the
\cpp{GroundStationManagerApp}, which aggregates them and recovers the orbit from
a perturbed a-priori state.

The tracking apps add the dominant Earth signal-path delays to the \emph{truth}
observable --- Saastamoinen troposphere ($\approx\SI{2.4}{\meter}$ at zenith,
$\approx\SI{14}{\meter}$ at the \SI{10}{\degree} mask), a thin-shell ionosphere
(sub-metre at X-band), the relativistic Shapiro delay (a few metres) and a solid
Earth tide station displacement of up to \SI{0.3}{\meter}. The manager removes
the deterministic terms in full (\code{model\_shapiro},
\code{model\_solid\_earth\_tide}) and calibrates a configurable fraction of the
atmospheric ones (\code{troposphere\_cancel\_fraction} and
\code{ionosphere\_cancel\_fraction}, here 0.8), inflating the range noise for
the uncancelled remainder via \code{residual\_delay\_noise\_scale}. An optional
\code{filter\_dynamics:} block lets the estimator run at a deliberately
different fidelity from truth.

\paragraph{The computation.}
The manager runs an iterative weighted-least-squares batch filter whose design
matrix is chained analytically through the autodiff state-transition matrix,
then a square-root information filter and a fixed-interval smoother over the
same data. The \SI{36}{\hour} arc is compute-heavy, so the notebook reads back
CSVs written by \file{python/examples/ex7_run_odts.py}.

\begin{lstlisting}[language=py, caption={Reading back the exported ODTS run.}, label={lst:tut-ex7-load}]
# python/examples/ex7_groundstation_odts.ipynb
DATA = seed_example_output("ex7_data")
meta = json.load(open(DATA / "meta.json"))
grid = pd.read_csv(DATA / "grid.csv")
iters = pd.read_csv(DATA / "iterations.csv")
elevation = pd.read_csv(DATA / "elevation.csv")
measurements = pd.read_csv(DATA / "measurements.csv")
epoch_sol = pd.read_csv(DATA / "epoch_solution.csv").set_index("which")
covariance = pd.read_csv(DATA / "covariance.csv", index_col=0).values
# duration=36:00:00  epochs=433  measurements=627  converged=True  iterations=5
\end{lstlisting}

\begin{lstlisting}[language=py, caption={Rotating the estimate error and its covariance into the RTN frame.}, label={lst:tut-ex7-rtn}]
# python/examples/ex7_groundstation_odts.ipynb
R_hat = r_true / np.linalg.norm(r_true, axis=1, keepdims=True)
h = np.cross(r_true, v_true)
N_hat = h / np.linalg.norm(h, axis=1, keepdims=True)
T_hat = np.cross(N_hat, R_hat)
Q = np.stack([R_hat, T_hat, N_hat], axis=1)

dpos_rtn = np.einsum("nij,nj->ni", Q, est[:, :3] - truth[:, :3])
Prr = np.einsum("nij,njk,nlk->nil", Q, est_cov[:, :3, :3], Q)
sig_pos_rtn = np.sqrt(np.diagonal(Prr, axis1=1, axis2=2))
\end{lstlisting}

The corresponding C++ driver reads the same solution off the manager
application:

\begin{lstlisting}[language=cpp, caption={Reading the centralized OD solution off the manager app.}, label={lst:tut-ex7-cpp}]
// cpp/examples/tutorials/ex7_groundstation_odts.cc
Simulation sim(YAML::LoadFile("configs/ground_station_odts.yaml"));
sim.Run();

auto mgr = std::dynamic_pointer_cast<GroundStationManagerApp>(
    sim.GetAgent("gs_manager")->GetApplication());
std::cout << "Batch filter converged: " << (mgr->Converged() ? "yes" : "no")
          << " in " << mgr->NumIterations() << " iterations\n";
const Vec6d x0_est = mgr->X0Estimated();
const double sigma_pos = std::sqrt(mgr->Covariance().topLeftCorner(3, 3).trace());
\end{lstlisting}

\paragraph{The results.}
Six figures, plus a convergence table.

Figure~1 (\cref{fig:tut-ex7-vis}) plots each station's topocentric elevation of
the satellite against time in hours, with the region above the \SI{10}{\degree}
mask shaded --- those shaded bands are exactly the passes that produced
measurements. The coverage is markedly uneven: Goldstone sees the satellite at
291 of 432 epochs (\SI{67.4}{\percent}), Madrid 221 (\SI{51.2}{\percent}) and
Canberra only 115 (\SI{26.6}{\percent}). Figure~2 shows the resulting
observations themselves: range in kilometres (a few hundred thousand, as
expected at lunar distance) and range rate in \si{\meter\per\second}, one colour
per station, visibly segmented into passes.

\begin{figure}[htbp]
  \centering
  \datafig[width=0.92\textwidth]{tutorials/ex7_groundstation_odts_1.png}
  \caption{Example~7, first figure: DSN elevation of the lunar satellite from
           Goldstone, Canberra and Madrid over \SI{36}{\hour}, with the
           above-mask tracking passes shaded.}
  \label{fig:tut-ex7-vis}
\end{figure}

Figure~3 (\cref{fig:tut-ex7-conv}) is the batch convergence, in two
logarithmic panels: the epoch-state error against iteration on the left, and the
RMS measurement residual against iteration on the right with dotted lines at the
assumed noise floors. The printed table is worth quoting because it shows a real
Gauss--Newton trajectory rather than a monotone one. Starting from an a-priori
that is \SI{1149.0}{\meter} and \SI{0.468}{\meter\per\second} from truth, the
range residual falls from \SI{116286.7}{\meter} to \SI{12811.7}{\meter} to
\SI{1791.4}{\meter} to \SI{10.5}{\meter} to \SI{10.2}{\meter}; the state error
\emph{worsens} at iteration~1 (\SI{16144.5}{\meter}) before collapsing to
\SI{24.4}{\meter}, \SI{3.96}{\meter} and finally \SI{1.639}{\meter}. The filter
converges in five iterations, the residual settles at the noise floor, and the
weighted RMS lands at 0.989 --- essentially unity, the signature of a correctly
weighted least-squares solution.

\begin{figure}[htbp]
  \centering
  \datafig[width=0.92\textwidth]{tutorials/ex7_groundstation_odts_3.png}
  \caption{Example~7, third figure: batch least-squares convergence. Left, the
           epoch-state error; right, the RMS range and range-rate residuals
           against their noise floors (dotted).}
  \label{fig:tut-ex7-conv}
\end{figure}

The final solution is a position error of \SI{1.639}{\meter} against a formal
$1\sigma$ of \SI{1.877}{\meter}, and a velocity error of
\SI{0.711}{\milli\meter\per\second} against a formal
\SI{0.973}{\milli\meter\per\second} --- error and covariance agree, so the
formal statistics are believable at the epoch. Figure~4 shows that solution
propagated: six panels of RTN position and velocity error along the arc inside
the batch covariance tube $\mat{P}(t_i)=\mat\Phi\mat{P}_0\mat\Phi\transpose$.

Figure~5 (\cref{fig:tut-ex7-srif}) is the methodological punchline. It plots
the formal position $1\sigma$ on a logarithmic axis for three estimators over
the same arc, with the tracking passes shaded: the batch
$\mat\Phi\mat{P}_0\mat\Phi\transpose$ curve, the SRIF forward filter, and the
SRIF fixed-interval smoother. The batch curve is optimistic --- it assumes
perfect dynamics, so mapping the epoch covariance forward never grows it
honestly. The forward filter, which carries white-noise-acceleration process
noise, grows visibly through each untracked apoapsis coast and is pulled back
down at every pass, giving the characteristic scalloped shape. The smoother,
which uses data on both sides of every epoch, stays bounded throughout. Figure~6
then shows the smoothed error in RTN with the smoother covariance tube (solid)
against the batch tube (dashed), confirming that the smoother's tighter
\emph{and} honest envelope still contains the error.

\begin{figure}[htbp]
  \centering
  \datafig[width=0.92\textwidth]{tutorials/ex7_groundstation_odts_5.png}
  \caption{Example~7, fifth figure: formal position $1\sigma$ from the batch
           mapping, the SRIF forward filter, and the SRIF smoother. Shaded
           bands are tracking passes; the forward filter grows through every
           untracked coast.}
  \label{fig:tut-ex7-srif}
\end{figure}

The batch and square-root information filters are specified in
\cref{chap:filters}; the closed forms and magnitudes of the ground-station
signal-path and station-location corrections are in \cref{chap:measurements}.

% ===========================================================================
\section{Example 8: Distributed Inter-Satellite-Link ODTS}
\label{sec:tut-ex8}

\textbf{Python:} \file{python/examples/ex8_isl_odts.ipynb} \quad
\textbf{C++:} \file{cpp/examples/tutorials/ex8_isl_odts.cc}

\paragraph{The scenario.}
NASA's five-satellite \textbf{LCRNS Reference Constellation~3.1}, taken as
Cartesian states in the Moon-centred ICRF at the reference epoch
\textbf{2027-03-01 00:00:00}, propagated over a \textbf{\SI{6}{\hour}} arc at a
\textbf{\SI{60}{\second}} measurement interval (361 epochs). The five ELFOs share
an orbital period of about \SI{30}{\hour}; their initial radii are
\SI{3485.5}{}, \SI{19149.8}{}, \SI{12548.6}{}, \SI{17631.0}{} and
\SI{17602.3}{\kilo\meter}, with \code{SV-1} and \code{SV-2} in one plane,
\code{SV-3} and \code{SV-4} in a second, and \code{SV-5} in a third. The
constellation is fully cross-linked --- four links per satellite --- and is
supplemented by three lunar surface stations.

The agent roster the notebook prints from \file{configs/isl_odts_distributed.yaml}
is five \cpp{Spacecraft} each hosting a \cpp{SatelliteOdtsApp}, three
\cpp{SurfaceStation} each hosting a \cpp{StationBeaconSensor}, and one
\cpp{SurfaceStationManager} hosting the centralized \cpp{GroundOdtsApp}. The
shared \code{world:} block holds the higher-fidelity truth force model; each
onboard app carries only its \emph{filter} model, a reduced $8\times8$ lunar
field. The two prior uncertainties are tuned deliberately: each satellite knows
its \emph{own} freshly initialised state to \SI{200}{\meter} but has only a
coarse broadcast ephemeris of its neighbours, \SI{500}{\meter}. Crosslink range
noise is \SI{1.0}{\meter} and range-rate noise
\SI{0.001}{\meter\per\second}; two-way time and frequency transfer are enabled;
consider states are exchanged every \SI{10}{\minute}.

\paragraph{The computation.}
Every satellite runs its own Schmidt EKF: it estimates its own eight-state
$[\vec{r},\vec{v},b,\dot b]$ and carries each neighbour as a \emph{consider}
state --- propagated and used in the measurement update, but never directly
corrected. Satellite-to-satellite links give two-way range and range rate (nearly
clock-free, so they constrain relative geometry) plus two-way time transfer
$c\,(b_i-b_j)$ and frequency transfer $c\,(d_i-d_j)$. Surface stations exchange a
one-way pseudorange and its Doppler with every satellite above their mask,
supplying the one \emph{absolute} time reference.

Five variants are produced by \file{python/examples/ex8_run_isl_odts.py} over
identical truth and geometry: the full system, exchange off, stations off, a
covariance-intersection exchange, and a three-step sweep of the crosslink clock
observables.

\begin{lstlisting}[language=py, caption={The LCRNS reference states used as truth initial conditions.}, label={lst:tut-ex8-lcrns}]
# python/examples/ex8_isl_odts.ipynb
t0_str = "2027-03-01T00:00:00"    # UTC, LCRNS Reference Constellation 3.1 epoch
lcrns_states_km = {   # Moon-centered ICRF, [km], [km/s]  (Ryden & Volle, Table 2)
    "SV-1": ([-198.931445, 386.023132, 3458.355412], [-1.539867, 0.022243, -0.091059]),
    "SV-2": ([1089.980598, -2260.918271, -18984.562823], [0.280301, -0.004049, 0.016575]),
    "SV-3": ([9187.665852, -260.470721, -8543.222759], [0.273629, 0.417588, -0.313828]),
    "SV-4": ([6484.019002, 14329.883921, -7966.345594], [-0.258433, 0.033951, 0.235264]),
    "SV-5": ([-5074.242314, 14473.022751, -8638.530033], [-0.229931, -0.026926, -0.264381]),
}
\end{lstlisting}

\begin{lstlisting}[language=py, caption={Per-epoch NEES, the covariance-consistency metric.}, label={lst:tut-ex8-nees}]
# python/examples/ex8_isl_odts.ipynb
def nees_ts(r, dof=3):
    """Per-epoch NEES averaged over the constellation.

    dof=3 -> position only, dof=8 -> full [r, v, clock_bias, clock_drift] state.
    A consistent filter sits near `dof`; NEES >> dof means the covariance is optimistic.
    """
    N, ns = len(r.t_s), len(r.satellite_names)
    out = np.full(N, np.nan)
    for k in range(N):
        vals = []
        for j in range(ns):
            e = own_block(r, j)[k, :dof] - r.truth_states[j][k, :dof]
            P = own_cov_8x8(r, j)[k, :dof, :dof]
            vals.append(float(e @ np.linalg.solve(P, e)))
        out[k] = np.mean(vals)
    return out
\end{lstlisting}

\paragraph{The results.}
Six figures, each answering a separate question.

Figure~1 (\cref{fig:tut-ex8-mesh}) is the geometry: the five ELFO truth tracks
over the arc, one colour per satellite, with the full crosslink mesh drawn as
grey segments at a single snapshot epoch. It makes plain why the constellation
works --- the satellites spend most of their long period near apolune over the
southern hemisphere, so a south-polar station sees most of them most of the
time. Figure~2 shows the observable geometry the onboard filters range across:
two-way crosslink range and range rate for the four links referenced to
\code{SV-1}; the \code{SV-1}--\code{SV-2} range alone spans \SI{14859}{} to
\SI{22635}{\kilo\meter} over the six hours.

\begin{figure}[htbp]
  \centering
  \datafig[width=0.55\textwidth]{tutorials/ex8_isl_odts_1.png}
  \caption{Example~8, first figure: the five LCRNS ELFOs over the \SI{6}{\hour}
           arc, with the full inter-satellite crosslink mesh at one snapshot
           epoch.}
  \label{fig:tut-ex8-mesh}
\end{figure}

Figure~3 (\cref{fig:tut-ex8-conv}) is the convergence panel. The top row plots
each satellite's own position error and clock-bias error against time on
logarithmic axes with the individual $3\sigma$ envelopes dashed; the bottom row
isolates the two observability contributions using the constellation mean. The
numbers underneath separate them cleanly. Final own-state position error with the
full system is \SI{31.0}{}, \SI{25.4}{}, \SI{39.0}{}, \SI{55.7}{} and
\SI{53.0}{\meter} for \code{SV-1}--\code{SV-5}, a constellation mean of
\SI{40.8}{\meter}; with the \SI{10}{\minute} consider-state exchange
\emph{disabled} the mean degrades to \SI{162.6}{\meter}, with individual
satellites as bad as \SI{444}{\meter}. Clock bias is the mirror image: the mean
is \SI{24.9}{\meter} with the surface stations on and \SI{410.2}{\meter} with
them off, essentially unobserved. \textbf{Crosslinks plus exchange fix the
orbit; the surface station fixes the time.}

\begin{figure}[htbp]
  \centering
  \datafig[width=\textwidth]{tutorials/ex8_isl_odts_3.png}
  \caption{Example~8, third figure: own position and clock-bias error for all
           five onboard filters with their $3\sigma$ envelopes (top), and the
           exchange-off and station-off ablations against the constellation mean
           (bottom).}
  \label{fig:tut-ex8-conv}
\end{figure}

Figure~4 asks whether the reported covariance is \emph{honest}, using the
normalised estimation error squared, which for a consistent filter sits near its
degrees of freedom (3 for position). The naive-overwrite exchange --- each
filter simply adopting a neighbour's self-reported covariance --- reaches a
position NEES of \num{34180}, about \num{11393} times the consistent value of 3:
it reports a $\sigma$ of \SI{0.42}{\meter} while the true RMS error is
\SI{24.5}{\meter}. The cause is named explicitly: the five filters share the same
crosslink observations and have already been mixed by previous exchanges, so
treating a broadcast as fresh independent information counts it repeatedly, a
data-incest feedback loop. Switching to a covariance-intersection exchange with
matched process noise brings the NEES down to \num{236} --- a $145\times$
reduction in over-confidence --- at the cost of only a slight accuracy change
(RMS \SI{19.5}{\meter}) and a reported $\sigma$ of \SI{16.25}{\meter}, now the
same order as the true error.

Figure~5 (\cref{fig:tut-ex8-transfer}) sweeps the crosslink clock observables
on the same truth with the consistent exchange: constellation-mean clock-bias
error and clock-drift error against time for three configurations. Range and
range rate alone leave \SI{58.5}{\meter} of clock bias and
\SI{0.00356}{\meter\per\second} of drift; adding two-way time transfer collapses
the bias to \SI{19.0}{\meter} and the drift to
\SI{0.00121}{\meter\per\second}; adding frequency transfer as well gives
\SI{5.4}{\meter} and \SI{0.00038}{\meter\per\second}. Coherent two-way ranging
cancels the clocks and therefore says nothing about them; the two transfer
observables are what make the constellation's relative clocks observable.

\begin{figure}[htbp]
  \centering
  \datafig[width=\textwidth]{tutorials/ex8_isl_odts_5.png}
  \caption{Example~8, fifth figure: constellation-mean clock-bias and
           clock-drift error as two-way time transfer and then frequency
           transfer are added to the crosslink observables.}
  \label{fig:tut-ex8-transfer}
\end{figure}

Figure~6 (\cref{fig:tut-ex8-central}) is the head-to-head against the classic
architecture: per-satellite final position error, clock error, and a summary
panel, distributed ISL against the centralized ground filter that sees only the
surface-station pseudoranges. The distributed scheme is \emph{uniform} ---
\SIrange{25}{56}{\meter} across all five satellites --- while the ground filter
ranges from \SI{1503}{\meter} on the well-covered \code{SV-5} to
\SI{6789}{\meter} on the low-perilune \code{SV-1} that spends long stretches
below every station's horizon, with clock errors from \SI{1455}{} to
\SI{14371}{\meter}. Crosslinks buy observability that no ground network alone
can supply.

\begin{figure}[htbp]
  \centering
  \datafig[width=\textwidth]{tutorials/ex8_isl_odts_6.png}
  \caption{Example~8, sixth figure: final per-satellite position and clock error,
           distributed inter-satellite-link filters versus the centralized
           ground filter on identical truth and station geometry.}
  \label{fig:tut-ex8-central}
\end{figure}

The Schmidt/consider formulation, the covariance-intersection fusion, the
stochastic-cloning bookkeeping, and the covariance conventions are specified in
\cref{chap:filters}; the crosslink and beacon observables are in
\cref{chap:measurements}.

% ===========================================================================
\section{Example 9: Ephemeris and Almanac Fitting}
\label{sec:tut-ex9}

\textbf{Python:} \file{python/examples/ex9_ephemeris.ipynb} \quad
\textbf{C++:} \file{cpp/examples/tutorials/ex9_ephemeris.cc}

\paragraph{The scenario.}
The accuracy-versus-broadcast-datasize trade-off of two lunar-satellite orbit
models fit to a numerically propagated truth trajectory. The truth is the ELFO
of \cref{sec:tut-ex1} propagated under N-body dynamics and exported already
rotated into the Moon-fixed principal-axis frame \code{MOON\_PA} --- the frame a
lunar navigation broadcast uses, so that a surface user can form topocentric
geometry directly. Two arcs are exported: a dense \textbf{1-day arc, 1441
samples at \SI{60}{\second}}, for the precise ephemeris, and a coarse
\textbf{30-day arc, 4321 samples at \SI{600}{\second}}, for the almanac.

The two models are \cpp{LansEphemeris} --- an osculating two-body baseline plus
per-axis Chebyshev Cartesian residuals, fit here over a \textbf{2-hour} window
with polynomial order~8, giving \textbf{35 parameters} --- and
\cpp{LansAlmanac}, which fits each element with a linear polynomial plus one
element-specific Fourier term, replacing the argument of periapsis by an
argument-of-latitude correction, fit over a \textbf{15-day} window with
\textbf{27 parameters}.

\begin{lstlisting}[language=py, caption={Fitting both broadcast models in the Moon-fixed frame.}, label={lst:tut-ex9-fit}]
# python/examples/ex9_ephemeris.ipynb
cart_opts = pnt.EphemerisFitOptions()
cart_opts.poly_order = 8
cart_opts.use_keplerian_baseline = True
cart_opts.frame = pnt.MOON_PA
cart_eph = pnt.LansEphemeris(cart_opts)

alm_opts = pnt.AlmanacFitOptions()
alm_opts.poly_order = 1        # linear polynomial + one Fourier term per element
almanac = pnt.LansAlmanac(alm_opts)
\end{lstlisting}

\paragraph{The computation.}
Both models are fit and evaluated in \code{MOON\_PA}; internally the osculating
elements are formed in the quasi-inertial principal-axis-inertial (PAI) frame,
which shares position with PA and differs in velocity by
$\vec\omega\times\vec{r}$ with $\omega_M=\SI{2.6617e-6}{\radian\per\second}$,
and the two-body baseline's node drifts at the Moon's rotation rate so that the
model lands in the rotating output frame.

The trade-off study is automated by an \cpp{EphemerisApp} hosted on a thin
\cpp{EphemerisManager} agent, run out of process by
\file{python/examples/ex9_run_ephemeris.py}. For each fit-window length it fits
several windows, and for each parameter a bisection finds the fewest fractional
bits whose quantisation keeps the \emph{worst-case} in-window position offset
under the target precision; the broadcast value is then quantised to that
resolution and the orbit re-evaluated, so the reported accuracy is what a
receiver decoding the message actually gets.

\begin{lstlisting}[language=py, caption={Per-parameter bit search, mirroring the C++ EphemerisApp.}, label={lst:tut-ex9-bits}]
# python/examples/ex9_ephemeris.ipynb
def _search_param_bits(eph_model, t_s, params, idx, precision_m, max_bits=50):
    base = eph_model.eval(t_s, params)[:, :3]

    def perturbed_err(k):
        eps = 2.0 ** (-k)
        q = params.copy()
        q[idx] += eps
        plus = eph_model.eval(t_s, q)[:, :3]
        q[idx] -= 2.0 * eps
        minus = eph_model.eval(t_s, q)[:, :3]
        return max(...)
\end{lstlisting}

\begin{lstlisting}[language=py, caption={Turning the almanac into the observables a surface user needs.}, label={lst:tut-ex9-obs}]
# python/examples/ex9_ephemeris.ipynb
user_lla = np.array([-90.0, 0.0, 0.0])            # lunar South Pole [lat deg, lon deg, alt m]
user_pa = pnt.lat_lon_alt_to_cart(user_lla, pnt.R_MOON)
f_carrier_hz = 2.4e9                              # nominal S-band carrier

def observables(rv_pa):
    """(az[deg], el[deg], range[m], range-rate[m/s], Doppler[Hz]) at the user."""
    aer = np.asarray(pnt.cart_to_az_el_range(rv_pa[:, :3], user_pa))
    rel = rv_pa[:, :3] - user_pa
    range_rate = np.sum(rel * rv_pa[:, 3:], axis=1) / np.linalg.norm(rel, axis=1)
    doppler_hz = -range_rate * f_carrier_hz / pnt.C
    return np.rad2deg(aer[:, 0]), np.rad2deg(aer[:, 1]), aer[:, 2], range_rate, doppler_hz
\end{lstlisting}

\paragraph{The results.}
Four figures plus two tables, and the headline is the four-orders-of-magnitude
gap between the two models.

The single-window comparison prints, for the 2-hour \cpp{LansEphemeris},
position RMS \SI{0.183}{\meter} and 95th percentile \SI{0.318}{\meter}, with
velocity RMS \SI{1.54}{\milli\meter\per\second}; for the 15-day
\cpp{LansAlmanac}, position RMS \SI{11156}{\meter} and 95th percentile
\SI{19558}{\meter}, with velocity RMS \SI{1.87}{\meter\per\second}. Figure~1
(\cref{fig:tut-ex9-fit}) plots those two error histories side by side on
logarithmic ordinates --- hours within the window for the ephemeris, days for
the almanac. The ephemeris error is essentially flat and sub-metre across the
window; the almanac error oscillates strongly at the orbital period, spiking
near periapsis where the geometry changes fastest.

\begin{figure}[htbp]
  \centering
  \datafig[width=\textwidth]{tutorials/ex9_ephemeris_1.png}
  \caption{Example~9, first figure: 3-D position error inside one fit window,
           \cpp{LansEphemeris} over \SI{2}{\hour} (left) and \cpp{LansAlmanac}
           over \SI{15}{\dayunit} (right). Note the different ordinate decades.}
  \label{fig:tut-ex9-fit}
\end{figure}

Figure~2 (\cref{fig:tut-ex9-user}) asks what that almanac error means for the
user who has to point an antenna: four stacked panels of azimuth, elevation,
range and Doppler prediction error for a receiver at the lunar South Pole over
the 15-day window, with the intervals in which the satellite is below a
\SI{5}{\degree} elevation mask greyed out. Over the 1625 of 2161 samples
(\SI{75.2}{\percent}) in which the satellite is actually visible, the errors are
azimuth RMS \SI{0.082}{\degree} (95th percentile \SI{0.160}{\degree}, max
\SI{0.242}{\degree}), elevation RMS \SI{0.049}{\degree}, range RMS
\SI{4036}{\meter} (max \SI{13640}{\meter}) and Doppler RMS \SI{3.12}{\hertz}
at S-band (95th percentile \SI{6.25}{\hertz}). This is the practical reading of
``coarse but useful'': kilometres of range error, but pointing good to a tenth
of a degree and Doppler good to a few hertz --- enough to acquire the signal,
which is all an almanac has to do.

\begin{figure}[htbp]
  \centering
  \datafig[width=0.78\textwidth]{tutorials/ex9_ephemeris_2.png}
  \caption{Example~9, second figure: \cpp{LansAlmanac} prediction error in
           azimuth, elevation, range and S-band Doppler as seen from the lunar
           South Pole. Grey bands are below the \SI{5}{\degree} elevation mask.}
  \label{fig:tut-ex9-user}
\end{figure}

Figure~3 (\cref{fig:tut-ex9-trade}) is the trade-off proper: 95th-percentile
3-D position error against fit-window length on the left, and against total
broadcast size in bits on the right, with each point annotated by its window.
The six sweep points are, for \cpp{LansEphemeris}, \SI{30}{\minute} / 29
parameters / \textbf{381 bits} / \SI{0.048}{\meter}, \SI{60}{\minute} /
\textbf{400 bits} / \SI{0.102}{\meter}, and \SI{120}{\minute} / \textbf{428
bits} / \SI{0.345}{\meter}; and for \cpp{LansAlmanac}, \SI{1}{\dayunit} / 27
parameters / \textbf{740 bits} / \SI{8452}{\meter}, \SI{7}{\dayunit} /
\textbf{752 bits} / \SI{14141}{\meter}, and \SI{15}{\dayunit} / \textbf{754
bits} / \SI{26889}{\meter}. Two conclusions fall out. First, within a model,
error grows steeply with validity while the bit count barely moves --- the
almanac buys fifteen days for fourteen extra bits, at the price of a factor of
three in error. Second, the two models are not on one curve: the almanac costs
roughly \emph{twice} the bits of the ephemeris and delivers four orders of
magnitude worse accuracy, because it is answering a different question
(acquisition over weeks, not ranging over hours).

\begin{figure}[htbp]
  \centering
  \datafig[width=\textwidth]{tutorials/ex9_ephemeris_3.png}
  \caption{Example~9, third figure: 95th-percentile position error against fit
           window (left) and against broadcast size in bits (right) for both
           models, each point labelled with its window length.}
  \label{fig:tut-ex9-trade}
\end{figure}

Figure~4 opens up the 428-bit ephemeris point into a per-parameter bar chart:
the six Keplerian baseline elements and the three per-axis Chebyshev blocks
stacked by coefficient order. The bars sum to exactly the 428 bits of the
trade-off point, so the reader can see which parameters are expensive --- the
low-order coefficients, which move the orbit most per unit change --- and which
are nearly free.

Because the sweep is compute-heavy, the notebook delegates it to
\file{python/examples/ex9_run_ephemeris.py} and reads back the exported results.
The parameterisations, the fitting procedure, the quantisation and bit-budget
accounting are specified in \cref{chap:ephemeris-almanac}.

% ===========================================================================
\section{Example 10: Surface Rover Navigation}
\label{sec:tut-ex10}

\textbf{Python:} \file{python/examples/ex10_surface_rover.ipynb} \quad
\textbf{C++:} \file{cpp/examples/tutorials/ex10_surface_rover.cc}

\paragraph{The scenario.}
A lunar-surface rover driving a slowly curving arc at \textbf{Site01
(``Connecting ridge'')} near the south pole, over \textbf{2401 epochs at
\SI{1}{\second}}, i.e.\ a \textbf{\SI{40}{\minute}} arc. It fuses a full
strapdown inertial measurement unit (accelerometer and gyroscope, Kalibr noise
model), LCRNS/LANS pseudoranges from a five-satellite relay constellation, and a
LOLA digital-elevation-model altitude constraint, estimating the IMU biases
online.

The scenario is \file{configs/surface_rover_nav.yaml}. A shared \code{World}
holds the frame (\code{MOON\_PA}), point-mass lunar gravity, and the DEM site,
downloaded and cached from NASA PGDA. A thin \code{Rover} agent hosts the
\cpp{SurfaceRoverNavApp}, an error-state strapdown-INS EKF that precomputes the
driving-arc truth, synthesises the IMU, pseudorange and DEM measurements each
epoch from the \code{World}, and runs predict/update. A
\code{LunarNavConstellation} agent stands in for the five relay satellites ---
one configuration entry expanding into five self-propagating \cpp{Spacecraft}
--- so the rover queries the relay geometry rather than propagating the orbits
itself.

The \SI{30}{\minute}-scale INS arc is compute-heavy, so
\file{python/examples/ex10_run_surface_nav.py} runs it once (a full-sensor run
and a DEM-disabled ablation) and caches CSVs that the notebook reads.

\begin{lstlisting}[language=py, caption={Reading the full-sensor run and the DEM ablation.}, label={lst:tut-ex10-load}]
# python/examples/ex10_surface_rover.ipynb
DATA = seed_example_output("ex10_data")
meta = json.load(open(DATA / "meta.json"))
on = pd.read_csv(DATA / "series_dem_on.csv")    # DEM constraint enabled
off = pd.read_csv(DATA / "series_dem_off.csv")  # DEM constraint disabled (ablation)
dem_x = np.load(DATA / "dem_x.npy")
dem_y = np.load(DATA / "dem_y.npy")
dem_z = np.load(DATA / "dem_elevation.npy")
# Site: Site01 (Connecting ridge)  epochs=2401  dt=1.0 s
\end{lstlisting}

The C++ twin performs the same ablation by editing the configuration in memory
rather than on disk:

\begin{lstlisting}[language=cpp, caption={Ablating the DEM constraint by editing the config in memory.}, label={lst:tut-ex10-cpp}]
// cpp/examples/tutorials/ex10_surface_rover.cc
Config cfg_off = YAML::LoadFile(config_path);
cfg_off["agents"]["Rover"]["application"]["enable_dem_constraint"] = false;
Simulation sim_off(cfg_off);
sim_off.Run();
auto* app_off = dynamic_cast<SurfaceRoverNavApp*>(
    sim_off.GetAgent("Rover")->GetApplication().get());
// app_off->pos_err_norm(), pos_err_enu(), n_visible(), site_id()
\end{lstlisting}

\paragraph{The results.}
Six figures, and one clean ablation number.

Figure~1 (\cref{fig:tut-ex10-track}) is the setting: a filled contour map of the
LOLA terrain elevation in kilometres over the local East/North tile, with the
truth ground track in white, the DEM-aided estimate dashed in red, and the start
marked. The two tracks are visually indistinguishable at map scale --- the
metre-level errors that follow are invisible against a multi-kilometre tile,
which is itself worth seeing.

\begin{figure}[htbp]
  \centering
  \datafig[width=0.66\textwidth]{tutorials/ex10_surface_rover_1.png}
  \caption{Example~10, first figure: LOLA terrain at Site01 with the rover's
           truth track (white) and DEM-aided estimate (red dashed).}
  \label{fig:tut-ex10-track}
\end{figure}

Figure~2 (\cref{fig:tut-ex10-err}) is the consistency check: East, North and Up
position error against time in minutes, each with the filter's $\pm3\sigma$
envelope, with the ordinate zoomed to the converged behaviour rather than the
initial transient. All three components stay inside the envelope, so the filter
is consistent in every axis. Figure~3 shows the online IMU bias estimation ---
accelerometer bias error in \si{\milli\meter\per\second\squared} and gyroscope
bias error in \si{\micro\degree\per\second}, per axis, against the dotted
$3\sigma$ bounds --- converging toward zero as the arc proceeds. Figure~4 shows
the attitude error (the nav-frame rotation vector, roll/pitch/yaw in degrees)
inside its own $3\sigma$ envelope; attitude observability here comes from the
accelerometer sensing gravity plus the rover's own motion.

\begin{figure}[htbp]
  \centering
  \datafig[width=0.86\textwidth]{tutorials/ex10_surface_rover_2.png}
  \caption{Example~10, second figure: East/North/Up position error against the
           filter's $\pm3\sigma$ envelope, DEM-aided.}
  \label{fig:tut-ex10-err}
\end{figure}

Figure~5 (\cref{fig:tut-ex10-dem}) is the point of the example: the absolute
vertical error with the DEM constraint on and off, on one axis. The final
$|\mathrm{Up}|$ error is \textbf{\SI{1.21}{\meter}} with the constraint and
\textbf{\SI{16.42}{\meter}} without it, and the total 3-D position error goes
from \SI{6.06}{\meter} to \SI{16.65}{\meter}. Essentially all of the degradation
is vertical: with relays clustered above a polar user the vertical channel is
the weakly observable one, and tying the estimate to the terrain height is what
holds it down.

\begin{figure}[htbp]
  \centering
  \datafig[width=0.86\textwidth]{tutorials/ex10_surface_rover_5.png}
  \caption{Example~10, fifth figure: the DEM altitude constraint ablation.
           Removing it inflates the vertical error from
           $\approx\SI{1}{\meter}$ to $\approx\SI{16}{\meter}$.}
  \label{fig:tut-ex10-dem}
\end{figure}

Figure~6 closes the loop on the aiding source: the number of LCRNS relays above
the elevation mask on top, and the receiver-clock-bias error in nanoseconds with
its $\pm3\sigma$ band below, so the reader can correlate dips in relay
visibility with excursions in the clock estimate.

The IMU error model, the pseudorange observable, and the DEM height constraint
are specified in \cref{chap:measurements}; the error-state filter is in
\cref{chap:filters}.

% ===========================================================================
\section{Example 11: Lunar Lander Navigation}
\label{sec:tut-ex11}

\textbf{Python:} \file{python/examples/ex11_lander_navigation.ipynb} \quad
\textbf{C++:} \file{cpp/examples/tutorials/ex11_lander_navigation.cc}

\paragraph{The scenario.}
A lander on powered descent to the same south-pole site (\textbf{Site01,
Connecting ridge}), over \textbf{601 epochs at \SI{0.5}{\second}} --- a
\textbf{\SI{300}{\second}} descent. It glides from roughly \SI{2.5}{\kilo\meter}
downrange and \SI{2}{\kilo\meter} altitude to a hover, ending at a truth height
above terrain of \SI{15.00}{\meter}.

Navigation is an error-state INS multiplicative EKF fusing a full IMU (Kalibr
noise model), a nadir radar altimeter measuring height above the DEM terrain,
crater-landmark bearings against a synthetic crater map (terrain-relative
navigation), and LunaNet/LANS pseudoranges, with the IMU biases estimated
online. The single \code{Lander} agent hosts \emph{two} applications: a
\cpp{LanderGncApp} that owns the powered-descent truth trajectory --- a built-in
smoothstep, or a guidance-law trajectory injected through
\py{set\_reference\_trajectory\_enu} --- and writes the lander truth state each
epoch, and a \cpp{LanderNavApp} that reads that truth to synthesise its
measurements and drive the descent MEKF. The scenario is
\file{configs/lander_nav.yaml}; \file{python/examples/ex11_run_lander_nav.py}
runs the full-sensor descent plus a sensor ablation and caches the CSVs.

\begin{lstlisting}[language=py, caption={The cached descent run, crater map, and ablation table.}, label={lst:tut-ex11-load}]
# python/examples/ex11_lander_navigation.ipynb
DATA = seed_example_output("ex11_data")
meta = json.load(open(DATA / "meta.json"))
s = pd.read_csv(DATA / "series_all.csv")
craters = pd.read_csv(DATA / "craters.csv")
ablation = pd.read_csv(DATA / "ablation.csv")
# Site: Site01 (Connecting ridge)  epochs=601  dt=0.5 s
# final 3D position error (all sensors): 5.06 m ; touchdown altitude (truth): 15.00 m
\end{lstlisting}

\begin{lstlisting}[language=cpp, caption={Two applications on one agent, and the sensor ablation.}, label={lst:tut-ex11-cpp}]
// cpp/examples/tutorials/ex11_lander_navigation.cc
Simulation sim(YAML::LoadFile(config_path));
sim.Run();
auto* app = dynamic_cast<LanderNavApp*>(
    sim.GetAgent("Lander")->GetApplicationByName("LanderNavApp").get());
// app->alt_truth(), app->pos_err_norm(), app->vel_err_norm()

Config c = YAML::LoadFile(config_path);
c["agents"]["Lander"]["application"]["enable_craters"] = false;   // then re-run
\end{lstlisting}

\paragraph{The results.}
Five figures.

Figure~1 (\cref{fig:tut-ex11-track}) draws the descent over the terrain: a
filled elevation contour with the truth ground track in white, the estimated
track dashed in red, the synthetic crater field as small black dots, the start
marked with a star and the landing target with a cross. It shows both that the
descent is a few kilometres of downrange travel and that the crater map the
terrain-relative navigation uses actually covers the approach corridor.

\begin{figure}[htbp]
  \centering
  \datafig[width=0.66\textwidth]{tutorials/ex11_lander_navigation_1.png}
  \caption{Example~11, first figure: powered-descent ground track over the
           Site01 terrain, with the synthetic crater map used for
           terrain-relative navigation.}
  \label{fig:tut-ex11-track}
\end{figure}

Figure~2 pairs the altitude profile --- truth against estimated height above
terrain, falling from about \SI{2}{\kilo\meter} to the \SI{15}{\meter} hover ---
with the per-epoch counts of craters in the field of view and LunaNet satellites
visible, so the reader can see which aiding source is actually available at each
phase of the descent. Figure~3 (\cref{fig:tut-ex11-err}) is the consistency
plot: East, North and Up position error against the filter's $\pm3\sigma$
envelope over the \SI{300}{\second} descent. Figure~4 stacks the accelerometer
bias error, gyroscope bias error and attitude error with their bounds.

\begin{figure}[htbp]
  \centering
  \datafig[width=0.86\textwidth]{tutorials/ex11_lander_navigation_3.png}
  \caption{Example~11, third figure: descent position error in East/North/Up
           against the MEKF's $\pm3\sigma$ envelope.}
  \label{fig:tut-ex11-err}
\end{figure}

Figure~5 (\cref{fig:tut-ex11-ablation}) is the sensor ablation, as a labelled
bar chart of final 3-D position error: \textbf{all sensors
\SI{5.06}{\meter}}, \textbf{no craters \SI{5.11}{\meter}}, \textbf{no altimeter
\SI{4.76}{\meter}}, \textbf{no LunaNet \SI{12.15}{\meter}}. The honest reading
is that in \emph{this} configuration only LunaNet is load-bearing for the final
error: removing it more than doubles it, while removing craters or the altimeter
changes it by a few percent --- the altimeter case even improves marginally,
which is within the noise of a single run. Craters constrain the horizontal
channel and the altimeter the vertical, but with an absolute pseudorange fix
available throughout, neither is what sets the terminal accuracy. That is a
useful negative result: it says the ablation must be re-run without LunaNet if
one wants to see what terrain-relative navigation is worth on its own.

\begin{figure}[htbp]
  \centering
  \datafig[width=0.62\textwidth]{tutorials/ex11_lander_navigation_5.png}
  \caption{Example~11, fifth figure: final 3-D position error with each aiding
           source disabled in turn. Only the loss of LunaNet moves the terminal
           error appreciably.}
  \label{fig:tut-ex11-ablation}
\end{figure}

The altimeter, bearing, and pseudorange observables and their Jacobians are
specified in \cref{chap:measurements}; the multiplicative error-state
formulation is in \cref{chap:filters}.

% ===========================================================================
\section{Example 12: Lunar Navigation Constellation Design}
\label{sec:tut-ex12}

\textbf{Python:} \file{python/examples/ex12_constellation_design.ipynb} \quad
\textbf{C++:} none (Python only)

\paragraph{The scenario.}
The first example that studies geometry rather than estimating a state. A
reusable \py{LunarNavConstellation} class generates a symmetric ELFO
Walker-type constellation, propagates it, computes visibility and dilution of
precision, estimates the required EIRP, and optimises satellite phasing.

The orbit design is closed analytically before any propagation. A resonant
semi-major axis is chosen so that an integer number of revolutions fits exactly
into one sidereal month (\SI{27.3217}{\dayunit}), making the whole geometry
repeat monthly: \textbf{$N_\mathrm{rev}=22$} gives a period of
\SI{29.805}{\hour} and $a=\SI{11265.75}{\kilo\meter}$. With $e=0.60$ and
$\omega=\SI{90}{\degree}$ the frozen inclination follows from requiring the
secular apsidal rate to vanish, and comes out at
\textbf{\SI{51.757}{\degree}}, giving a perilune altitude of
\SI{2769}{\kilo\meter} over the northern hemisphere and an apolune altitude of
\SI{16288}{\kilo\meter} over the south --- exactly the dwell a polar service
needs.

The constellation is \textbf{6 satellites in 3 planes} (2 per plane), propagated
over one sidereal month at \SI{30}{\minute} sampling (1312 steps). The service
volume is the south-polar cap below \SI{75}{\degree}S sampled at
$\SI{5}{\degree}\times\SI{30}{\degree}$, \textbf{37 user points}, under a
\SI{10}{\degree} elevation mask, and the availability requirement is
$\mathrm{PDOP} < 6$.

\begin{lstlisting}[language=py, caption={Resonant semi-major axis and frozen inclination.}, label={lst:tut-ex12-design}]
# python/examples/ex12_constellation_design.ipynb
def resonant_sma(n_rev, gm=GM_MOON, t_repeat=T_SIDEREAL_MONTH):
    """Semi-major axis [m] and period [s] for n_rev revolutions per repeat window."""
    period = t_repeat / n_rev
    a = (gm * (period / (2 * np.pi)) ** 2) ** (1.0 / 3.0)
    return a, period


def omega_dot(a, e, inc, gm=GM_MOON, n_third=OMEGA_EM, j2=J2_MOON, r_body=R_MOON):
    """Secular apsidal rate [rad/s] at omega=90deg: Moon-J2 + Earth third-body."""
    n = np.sqrt(gm / a**3)
    p = a * (1.0 - e**2)
    w_j2 = 0.75 * n * j2 * (r_body / p) ** 2 * (5 * np.cos(inc) ** 2 - 1.0)
    w_3b = (0.75 * (n_third**2 / n) / np.sqrt(1 - e**2)
            * (2 + 3 * e**2 - 5 * np.sin(inc) ** 2))
    return w_j2 + w_3b


def frozen_inclination(a, e, **kw):
    return brentq(lambda inc: omega_dot(a, e, inc, **kw), 5 * RAD, 89 * RAD)
\end{lstlisting}

\paragraph{The computation.}
The design metric is position dilution of precision. With $\mat{G}$ the matrix
of line-of-sight unit vectors to the visible satellites, augmented with a column
of ones for the receiver clock,
\begin{equation}
  \mat{Q} = \left(\mat{G}\transpose \mat{G}\right)^{-1},
  \qquad
  \mathrm{PDOP} = \sqrt{Q_{xx} + Q_{yy} + Q_{zz}},
  \label{eq:tut-pdop}
\end{equation}
evaluated in a vectorised sweep over all users and epochs, with fewer than four
visible satellites flagged as no fix.

Phasing is then optimised. Re-propagating for each candidate would be far too
slow, so the notebook exploits two facts of the two-body resonant design: the
inertial-to-fixed rotation does not depend on phasing and can be precomputed
once, and shifting a satellite's mean anomaly by $\Delta M$ is exactly a time
shift $\Delta t = \Delta M/n$ along its inertially fixed orbit, so each plane's
reference trajectory can be precomputed over an extended grid and any phase
offset realised as an index roll. Each objective evaluation reduces to array
rolls, an \py{einsum}, and a batched $4\times4$ inverse.

\begin{lstlisting}[language=py, caption={Vectorised PDOP over a stack of line-of-sight vectors.}, label={lst:tut-ex12-pdop}]
# python/examples/ex12_constellation_design.ipynb
def batched_pdop(u_enu, visible):
    """Vectorized PDOP over a (..., n_sat, 3) stack of unit LOS vectors + visibility.

    Geometry matrix rows are [e, n, u, 1]; PDOP = sqrt(trace of the 3x3 position
    block of (H^T W H)^-1) with 0/1 weights. Entries with fewer than 4 visible
    satellites are set to DOP_INF.
    """
    h = np.concatenate([u_enu, np.ones(u_enu.shape[:-1] + (1,))], axis=-1)
\end{lstlisting}

\begin{lstlisting}[language=py, caption={Optimising the per-plane phase offsets.}, label={lst:tut-ex12-opt}]
# python/examples/ex12_constellation_design.ipynb
def neg_availability(x):
    return -evaluator.availability(np.concatenate([[0.0], x]))

result = differential_evolution(neg_availability,
                                bounds=[(0, 2 * np.pi)] * (N_PLANES - 1),
                                maxiter=20, popsize=12, tol=1e-4, seed=0, polish=False)
best_phase = np.concatenate([[0.0], result.x])
\end{lstlisting}

\paragraph{The results.}
Ten figures, in a narrative that goes design $\rightarrow$ baseline $\rightarrow$
optimisation $\rightarrow$ verification.

Figure~1 is the service volume: the 37 south-polar user points on a polar
projection with rings at \SI{5}{}, \SI{10}{} and \SI{15}{\degree} from the pole.
Figure~2 (\cref{fig:tut-ex12-const}) is the six-satellite constellation drawn in
\code{MOON\_CI} with the Moon to scale and the South Pole marked; the three
planes and the strongly eccentric, apolune-southern geometry are both apparent.
Figure~3 shows the ground tracks of all six satellites over one sidereal month
in the Moon-fixed frame, with the \SI{75}{\degree}S service boundary marked ---
the tracks converge over the southern cap, which is the whole design intent.

\begin{figure}[htbp]
  \centering
  \datafig[width=0.58\textwidth]{tutorials/ex12_constellation_design_2.png}
  \caption{Example~12, second figure: the six-satellite, three-plane frozen-ELFO
           navigation constellation in \code{MOON\_CI}, with the South Pole
           marked in red.}
  \label{fig:tut-ex12-const}
\end{figure}

Figures~4 and~5 establish the baseline with naive phasing (all plane offsets
zero). Figure~4 plots, for a user at the pole itself, the number of visible
satellites (with a red line at the four needed for a 3-D-plus-clock fix) and the
PDOP against time in days. Figure~5 maps per-user availability over the cap on a
polar projection, coloured red-to-green. With naive phasing the availability is
only \textbf{\SI{43.6}{\percent}}, with a median of 5 visible satellites (range
3 to 6) and a median PDOP of 3.57 --- the geometry is often adequate and often
not.

Figure~6 sizes the payload: required transmit EIRP against slant range, coloured
by elevation, for a \SI{30}{\dBHz} target with a \SI{6}{\dBi} user antenna at
the LunaNet Augmented Forward Signal band. The requirement runs from
\SI{-3.3}{\dBW} (\SI{0.46}{\watt}) at best to \textbf{\SI{4.5}{\dBW}}
(\SI{2.8}{\watt}) at the maximum in-view slant range of
\SI{17140}{\kilo\meter}; the 95th percentile is \SI{4.3}{\dBW}. A few watts
closes the link, so the constellation is not power-limited.

Figure~7 (\cref{fig:tut-ex12-land}) is the optimisation landscape: availability
as a function of the two free per-plane phase offsets on a
$25\times25$ grid (about \SI{15}{\degree} resolution), with the grid optimum
starred. It spans \textbf{\SI{20.2}{\percent}} at the worst phasing to
\textbf{\SI{69.2}{\percent}} at the best --- phasing alone is worth a factor of
three in availability, at zero cost in mass, power or satellite count. The
differential-evolution refinement finds offsets of
$(0,\,209.3,\,328.0)\si{\degree}$ and \textbf{\SI{72.0}{\percent}}.

\begin{figure}[htbp]
  \centering
  \datafig[width=0.6\textwidth]{tutorials/ex12_constellation_design_7.png}
  \caption{Example~12, seventh figure: $\mathrm{PDOP}<6$ availability over the
           south-polar service volume as a function of the two free per-plane
           phase offsets. The landscape spans \SIrange{20}{69}{\percent}.}
  \label{fig:tut-ex12-land}
\end{figure}

Figure~8 (\cref{fig:tut-ex12-avail}) puts the naive and optimised per-user
availability maps side by side on a shared colour scale, and figure~9 shows the
PDOP cumulative distribution over the service volume for both, with a vertical
line at $\mathrm{PDOP}=6$: the optimised curve crosses that line far higher, and
its whole distribution is shifted left.

\begin{figure}[htbp]
  \centering
  \datafig[width=\textwidth]{tutorials/ex12_constellation_design_8.png}
  \caption{Example~12, eighth figure: per-user $\mathrm{PDOP}<6$ availability
           over the south-polar cap, naive phasing (left) versus optimised
           phasing (right).}
  \label{fig:tut-ex12-avail}
\end{figure}

The verification step is deliberately sobering. Re-propagating the optimised
constellation with full N-body dynamics (a $7\times1$ lunar field plus Earth and
Sun third bodies) over the same month drops availability from
\textbf{\SI{71.2}{\percent}} (two-body) to \textbf{\SI{60.3}{\percent}}: the slow
frozen libration reshapes the geometry away from the idealised two-body optimum,
so a production design would have to optimise directly on the perturbed
trajectories. Figure~10 confirms the orbits nevertheless remain frozen --- an
$e$--$\omega$ phase portrait librating around the frozen point, and a bounded
inclination history --- so the degradation is a phasing effect, not an orbit
instability.

There is no dedicated Part~\ref{part:math} chapter for constellation design; the
propagation it rests on is specified in \cref{chap:dynamics}, the line-of-sight
geometry and visibility gating in \cref{chap:measurements}, and the design study
itself is \citep{iiyama2025constellation}. The same optimiser is reused as a
building block by \cref{sec:tut-ex10}, whose \code{LunarNavConstellation} agent
plays the relay constellation.

% ===========================================================================
\section{Example 13: Visualising Constellations in Cesium}
\label{sec:tut-ex13}

\textbf{Python:} \file{python/examples/ex13_cesium.ipynb} \quad
\textbf{C++:} none (Python only)

\paragraph{The scenario.}
Turning \lupnt{} trajectory samples into interactive CesiumJS scenes. Beyond
presentation, this is a practical debugging tool for frame conventions,
body-fixed coordinates, surface locations, and moving link geometry. Three
scenes are built.

The Earth scene loads the \textbf{31 GPS and 31 Galileo} two-line elements
archived at \textbf{2025-01-01} and propagates each with two-body dynamics for
\textbf{\SI{16}{\hour} at \SI{4}{\minute}} steps --- more than one full
revolution for both systems --- handing the inertial states to
\py{add\_trajectory}, which rotates them into the Earth-fixed frame. The lunar
scene takes the same five LCRNS Reference Constellation~3.1 states as
\cref{sec:tut-ex8}, at \textbf{2027-03-01}, and propagates them for
\textbf{\SI{31}{\hour} at \SI{10}{\minute}} steps (about one full ELFO period)
under an $8\times8$ lunar field plus Earth and Sun third bodies, rendering them
in the Moon-fixed \code{MOON\_PA} frame together with five surface sites ---
Shackleton HQ (\SI{-89.9}{\degree}, \SI{0}{\degree}), Connecting Ridge HQ
(\SI{-89.45}{\degree}, \SI{222.8}{\degree}), de Gerlache HQ
(\SI{-88.5}{\degree}, \SI{290}{\degree}), plus a near-side and a far-side site.
The third scene puts both constellations in one Earth-centred inertial view.

\begin{lstlisting}[language=py, caption={Building a CZML scene from inertial samples.}, label={lst:tut-ex13}]
# python/examples/ex13_cesium.ipynb
earth = pnt.plot.CesiumScene(body="EARTH", name="Earth GNSS constellation",
                             epoch=EPOCH_EARTH, multiplier=120, texture=True)
for i, tle in enumerate(gps_tles):
    earth.add_trajectory(f"GPS-{i+1}", t_tdb, gnss_eci(tle),
                         frame_in=pnt.ECI, color=(0, 180, 255))   # cyan
for i, tle in enumerate(gal_tles):
    earth.add_trajectory(f"GAL-{i+1}", t_tdb, gnss_eci(tle),
                         frame_in=pnt.ECI, color=(255, 180, 0))   # amber
\end{lstlisting}

\paragraph{The results.}
This notebook produces no raster figures --- its outputs are standalone HTML
files containing CZML, Cesium's time-dynamic JSON format, which is why it
contributes nothing to the extracted figure set. What it \emph{does} print is a
useful cross-check on the lunar scene: the perilune and apolune altitudes of the
five LCRNS satellites, \SI{1728}{}/\SI{17406}{}, \SI{1704}{}/\SI{17433}{},
\SI{1890}{}/\SI{17325}{}, \SI{1802}{}/\SI{17289}{} and
\SI{1723}{}/\SI{17379}{\kilo\meter} respectively --- five nearly identical
frozen ELFOs, as the reference constellation intends.

The notebook is candid about one rendering artefact: the drawn orbit path is the
finite propagated arc, not a mathematically closed conic, so the top of an ELFO
can look clipped where the first and last samples fail to wrap into one another.
\py{pnt.plot.CesiumScene} needs no Cesium ion account; it embeds a textured
ellipsoid and loads only the CesiumJS library from a public CDN. Surface assets
are placed with \py{add\_station(lat=..., lon=...)} in degrees, on a sphere of
the body's radius in the same frame as the orbits, so stations and satellites
share one consistent scene. The frame rotations \py{add\_trajectory} applies are
specified in \cref{chap:frame-conversions}.

% ===========================================================================
\section{Example 14: Optical Navigation}
\label{sec:tut-ex14}

\textbf{Python:} \file{python/examples/ex14_opnav.ipynb} \quad
\textbf{C++:} none (Python only)

\paragraph{The scenario.}
A compact optical-navigation pipeline for a spacecraft in lunar orbit,
combining \lupnt{} propagation, synthetic image generation in Blender, image
processing, geometric measurement formation, and an EKF.

The truth orbit is circular at \textbf{\SI{2600}{\kilo\meter}} altitude ---
orbit radius \SI{4337.4}{\kilo\meter}, inclination \SI{55}{\degree}, right
ascension \SI{35}{\degree}, period \SI{7.12}{\hour} --- at the epoch
\textbf{2025-01-01T00:00:00}. A high orbit is chosen deliberately so the full
lunar disk fits inside a modest field of view and a simple circle fit can
recover the limb. \textbf{18 images} are rendered over \textbf{0.36 of one
orbit}, a cadence of \SI{9.0}{\minute}, with a \textbf{$256\times256$ pixel}
camera at a \textbf{\SI{50}{\degree}} field of view, nadir-pointed with attitude
assumed known. The Sun is at \SI{146.74e6}{\kilo\meter} from the Moon and the
solar phase angle over the sequence runs from \SI{28.1}{} to \SI{67.8}{\degree},
so the images are genuinely phase-lit rather than fully illuminated.

Rendering uses Blender's Python module (\code{bpy}~4.5.11, provided by
\code{pixi install}) with NASA SVS CGI Moon Kit textures --- the 2025 LROC WAC
colour map for surface colour and the LOLA-derived preview displacement map as a
bump map --- plus directional sunlight, a weak camera-side fill so the night-side
limb stays visible, and a sparse star background.

\paragraph{The computation.}
Each image is segmented by taking the border median as the sky level, thresholding
at eight robust standard deviations above it, keeping the largest connected
component, filling interior mare and crater holes, extracting the one-pixel
exterior boundary, and fitting a circle to it by linear least squares (with a
half-pixel outward correction, because the binary boundary lies just inside the
continuous limb).

For a spherical Moon of radius $R_M$, a fitted image radius $r_\mathrm{px}$ at
focal length $f_\mathrm{px}$ gives
\begin{equation}
  \alpha = \arctan\!\left(\frac{r_\mathrm{px}}{f_\mathrm{px}}\right),
  \qquad
  \rho = \frac{R_M}{\sin\alpha},
  \label{eq:tut-opnav}
\end{equation}
and the fitted disk centre gives the bearing to the Moon centre, so each image
becomes a full Moon-centred inertial position measurement.

\begin{lstlisting}[language=py, caption={From a fitted limb circle to a Moon-centred position fix.}, label={lst:tut-ex14-meas}]
# python/examples/ex14_opnav.ipynb
def detection_to_position_ci(det, q_cam_to_ci, cam):
    alpha = np.arctan(det["radius_px"] / cam["f_px"])
    rho = R_MOON / np.sin(alpha)
    u_cam = np.array([(det["xc"] - cam["cx"]) / cam["f_px"],
                      -(det["yc"] - cam["cy"]) / cam["f_px"],
                      1.0])
    u_cam /= np.linalg.norm(u_cam)
    u_ci = q_cam_to_ci @ u_cam
    return -rho * u_ci
\end{lstlisting}

\begin{lstlisting}[language=py, caption={The position-only EKF measurement update.}, label={lst:tut-ex14-ekf}]
# python/examples/ex14_opnav.ipynb
def ekf_update_position(x, p, z, sigma_pos=8.0e3):
    h = np.zeros((3, 6))
    h[:, :3] = np.eye(3)
    r = (sigma_pos**2) * np.eye(3)
    y = z - h @ x
    s = h @ p @ h.T + r
    k = np.linalg.solve(s, h @ p).T
    x = x + k @ y
    i_kh = np.eye(6) - k @ h
    p = i_kh @ p @ i_kh.T + k @ r @ k.T
    return x, p, y
\end{lstlisting}

The filter state is Cartesian $[\vec{r},\vec{v}]$; propagation is two-body lunar
dynamics with RK4 and a process-noise acceleration of
\SI{2e-7}{\meter\per\second\squared}; the measurement is position-only with an
assumed $\sigma=\SI{8}{\kilo\meter}$. It is started deliberately imperfect:
the first image fix perturbed by \SI{15}{\kilo\meter} per axis, the truth
velocity perturbed by \SI{8}{\meter\per\second} per axis, with a prior of
\SI{25}{\kilo\meter} and \SI{20}{\meter\per\second}.

\paragraph{The results.}
Three figures, and a result that is worth reading carefully because it is not a
success story.

Figure~1 (\cref{fig:tut-ex14-limb}) shows the first six rendered frames with the
fitted limb circle drawn in red and the fitted radius in the panel title. The
renders are visibly phase-lit --- part of the disk is in shadow --- yet the fit
still tracks the full silhouette, because the segmentation keys on
\emph{difference from the sky}, not on brightness. This is the figure that
demonstrates the image-processing stage works.

\begin{figure}[htbp]
  \centering
  \datafig[width=0.82\textwidth]{tutorials/ex14_opnav_1.png}
  \caption{Example~14, first figure: six of the eighteen Blender-rendered lunar
           disks, each with the fitted limb circle and its radius in pixels.}
  \label{fig:tut-ex14-limb}
\end{figure}

The scalars quantify it: the limb radius error has an RMS of
\textbf{\SI{1.054}{px}}, and the resulting raw optical position fixes have an
RMS error of \textbf{\SI{31.87}{\kilo\meter}}. That number follows directly from
\cref{eq:tut-opnav}: at this range one pixel of radius error maps to tens of
kilometres of range error, so a one-pixel circle fit is a tens-of-kilometres
measurement. Optical limb navigation is a coarse absolute sensor, and the
example makes that unavoidable.

Figure~2 (\cref{fig:tut-ex14-err}) plots position error against time in hours
for the raw image fixes and for the EKF estimate on the left, and the velocity
error on the right. The filter ends at a final position error of
\SI{47.55}{\kilo\meter} and a filtered position RMS of
\SI{39.39}{\kilo\meter}, with a final velocity error of
\SI{16.02}{\meter\per\second}. Those are \emph{worse} than the raw fix RMS of
\SI{31.87}{\kilo\meter}. The reason is visible in the setup rather than the
plot: the arc is only 18 measurements over a third of an orbit, the filter is
started \SI{15}{\kilo\meter} and \SI{8}{\meter\per\second} away from truth with
a loose prior, and the assumed measurement $\sigma$ of \SI{8}{\kilo\meter} is
optimistic against the \SI{32}{\kilo\meter} the fixes actually deliver, so the
filter under-weights its own dynamics. The honest conclusion is that a
short-arc, position-only EKF cannot improve on the fixes it is fed; the
velocity, which no single image observes, is recovered only through the dynamic
model tying the sequence together, and \SI{16}{\meter\per\second} is what that
buys over a third of an orbit.

\begin{figure}[htbp]
  \centering
  \datafig[width=0.92\textwidth]{tutorials/ex14_opnav_2.png}
  \caption{Example~14, second figure: position error of the raw image fixes and
           of the EKF estimate (left), and the velocity error recovered through
           the dynamics (right).}
  \label{fig:tut-ex14-err}
\end{figure}

Figure~3 is the 3-D summary: the truth orbit, the EKF estimate as a dashed
crimson curve, and the individual image fixes as black markers, around a
half-transparent Moon. It shows the fixes scattering about the truth track and
the filtered curve running smoothly through the middle of them.

The notebook is explicit about what has been kept simple so the mathematical
structure stays visible: a spherical shape model, known camera attitude, no lens
distortion, no terrain relief, and no illuminated-limb-only fitting. Those are
listed as the natural extensions. The filter is specified in
\cref{chap:filters} and the bearing-measurement machinery in
\cref{chap:measurements}.

% ===========================================================================
\section{Example 15: Mars PNT}
\label{sec:tut-ex15}

\textbf{Python:} \file{python/examples/ex15_marspnt.ipynb} \quad
\textbf{C++:} \file{cpp/examples/tutorials/ex15_marspnt.cc}

\paragraph{The scenario.}
Nothing in \lupnt{}'s core is Moon-specific, and this example demonstrates that
by building a Mars navigation constellation end to end at the epoch
\textbf{2035-01-01}. Mars is taken from \py{Body.Mars} with
$GM=\SI{4.282838e13}{\meter\cubed\per\second\squared}$ and
$R=\SI{3396.0}{\kilo\meter}$, and the shipped \textbf{Mars-50c} spherical-harmonic
field truncated to \textbf{$8\times8$}. The measured sidereal day recovered from
the frame rotation is \SI{24.6230}{\hour} against the true \SI{24.6229}{\hour}.

The constellation is a \textbf{Walker-delta 9/3/1}: 9 satellites in 3 planes,
circular, \textbf{\SI{15000}{\kilo\meter}} altitude ($a=\SI{18396}{\kilo\meter}$),
\textbf{\SI{45}{\degree}} inclination \emph{to Mars' equator}, giving a period of
\SI{21.04}{\hour} that is deliberately \emph{not} synchronous with the
\SI{24.62}{\hour} Martian day, so the ground tracks walk westward and sample
every longitude within a few sols. It is propagated for \textbf{two sols} (985
epochs at \SI{180}{\second}, about 2.3 orbital periods).

The inclination caveat is worth stating: \code{MARS\_CI} is ICRF-aligned, and
Mars' spin pole is tilted \textbf{\SI{37.13}{\degree}} from the ICRF pole, so
elements built directly in \code{MARS\_CI} would be measured against the wrong
equator. The notebook therefore constructs a Mars-equatorial inertial frame from
\lupnt{}'s own Mars orientation --- rotating the body-fixed $z$-axis into
\code{MARS\_CI} to get the pole, then completing an orthonormal triad --- and
designs there.

\begin{lstlisting}[language=py, caption={A Mars-equatorial design frame from LuPNT's own orientation model.}, label={lst:tut-ex15-frame}]
# python/examples/ex15_marspnt.ipynb
def mars_equatorial_to_ci(t):
    """Rotation from a Mars-equatorial inertial frame (z = Mars pole) to MARS_CI."""
    pole = np.asarray(pnt.convert_frame(t, np.array([0.0, 0.0, 1.0]), FIXED, CI, True))
    pole /= np.linalg.norm(pole)
    x = np.cross([0.0, 0.0, 1.0], pole)   # node of Mars equator on the ICRF equator
    x /= np.linalg.norm(x)
    y = np.cross(pole, x)
    return np.column_stack([x, y, pole])  # columns = equatorial axes in MARS_CI
\end{lstlisting}

\begin{lstlisting}[language=py, caption={Walker layout, propagation, and rotation into the body-fixed frame.}, label={lst:tut-ex15-walker}]
# python/examples/ex15_marspnt.ipynb
def walker_elements(T, P, Fp, a, e, inc):
    """Classical elements [a,e,i,RAAN,argp,M] for a Walker-delta constellation."""
    spp = T // P
    coes = []
    for p in range(P):
        raan = p * 2 * np.pi / P
        for s in range(spp):
            M = (s * 2 * np.pi / spp + p * 2 * np.pi * Fp / T) % (2 * np.pi)
            coes.append([a, e, inc, raan, 0.0, M])
    return np.array(coes)

for k in range(N_SAT):
    POS_CI[k] = np.asarray(dyn.propagate(RV0[k], TS))[:, :3]
    POS_FIXED[k] = np.asarray(pnt.convert_frame(TS, POS_CI[k], CI, FIXED, True))
\end{lstlisting}

\paragraph{The computation.}
The same $\mathrm{PDOP} = \sqrt{Q_{11}+Q_{22}+Q_{33}}$ of \cref{eq:tut-pdop} is
evaluated on a $\SI{5}{\degree}\times\SI{5}{\degree}$ latitude/longitude grid at
every epoch under a \SI{5}{\degree} elevation mask, and the median over the
two-sol window is mapped. A user-positioning simulation at \textbf{Jezero
Crater} (\SI{18.44}{\degree}N, \SI{77.45}{\degree}E) then forms a pseudorange to
every visible satellite with $\sigma=\SI{8}{\meter}$ and solves for position and
clock bias by iterated Gauss--Newton least squares, seeded at the planet centre.

\paragraph{The results.}
Five figures.

Figure~1 (\cref{fig:tut-ex15-acc}) is the acceleration budget on a
representative satellite, drawn as a horizontal log-scale bar chart of the
strongest decomposed terms. The printed ranking is monopole
\SI{1.266e-1}{\meter\per\second\squared}, $J_2$
\SI{1.132e-5}{}, $C_{22}$ \SI{1.485e-6}{}, $C_{33}$
\SI{1.932e-7}{}, $C_{31}$ \SI{9.849e-8}{}, $C_{32}$
\SI{6.208e-8}{}, $J_3$ \SI{4.490e-8}{}, $C_{44}$
\SI{1.445e-8}{\meter\per\second\squared}. The point made in the text is that
$J_2$ --- about twice Earth's, at $1.96\times10^{-3}$ --- sits four orders below
the central term but acts \emph{secularly}, steadily rotating each orbit plane,
so it dominates the long-term geometry even though the instantaneous number
looks small. Propagation over two sols produces only \SI{6.9}{\kilo\meter} of
orbit-radius spread, the short-period $J_2$ breathing.

\begin{figure}[htbp]
  \centering
  \datafig[width=0.72\textwidth]{tutorials/ex15_marspnt_1.png}
  \caption{Example~15, first figure: Mars gravity acceleration budget on a
           satellite at $R_\mathrm{Mars}+\SI{15000}{\kilo\meter}$, decomposed
           per spherical-harmonic term.}
  \label{fig:tut-ex15-acc}
\end{figure}

Figure~2 is the constellation in 3-D over one period in \code{MARS\_CI},
coloured by plane over a textured Mars. Figure~3 draws the ground tracks over
two sols on the Mars surface map: each track is a slowly drifting sinusoid
reaching $\pm\SI{45}{\degree}$ latitude, with the three planes offset by
\SI{120}{\degree} in longitude, and the westward walk of the non-synchronous
period is directly visible.

Figure~4 (\cref{fig:tut-ex15-pdop}) is the coverage assessment: median PDOP over
the two sols on the left with a black contour at $\mathrm{PDOP}=6$ and Jezero
starred, and mean number of satellites in view on the right. The verdict is
blunt: the \textbf{global median PDOP is 6.74}, only \textbf{\SI{37}{\percent}}
of the grid achieves a median PDOP below 6, and the grid-average number of
satellites in view is \textbf{3.27} --- below the four needed for a fix. Nine
satellites is a genuinely minimal fleet, and the map shows exactly where the
gaps are: the mid-latitude band that the \SI{45}{\degree} inclination favours is
served, the poles are not.

\begin{figure}[htbp]
  \centering
  \datafig[width=\textwidth]{tutorials/ex15_marspnt_4.png}
  \caption{Example~15, fourth figure: median PDOP over two sols (left, black
           contour at PDOP 6) and mean satellites above the \SI{5}{\degree} mask
           (right), with Jezero Crater marked.}
  \label{fig:tut-ex15-pdop}
\end{figure}

Figure~5 (\cref{fig:tut-ex15-jezero}) is the user's experience at Jezero, in
three panels: 3-D position error against time coloured by PDOP, the error
cumulative distribution with the median marked, and a sky view at the
best-geometry epoch showing the visible satellites in azimuth and elevation
coloured by plane. A fix is available at \textbf{504 of 985 epochs}
(\SI{51}{\percent} of the window); when it is, the \textbf{median 3-D error is
\SI{30.0}{\meter}} with a 95th percentile of \SI{201.6}{\meter}
(\SI{15.2}{\meter} horizontal, \SI{20.0}{\meter} vertical), at a median PDOP of
4.84. Tens of metres from \SI{8}{\meter} ranging is exactly what a PDOP near 5
predicts, so the two halves of the analysis agree; the binding constraint is
availability, not accuracy.

\begin{figure}[htbp]
  \centering
  \datafig[width=\textwidth]{tutorials/ex15_marspnt_5.png}
  \caption{Example~15, fifth figure: pseudorange least-squares positioning at
           Jezero Crater --- error against time coloured by PDOP, the error CDF,
           and the sky view at the best-geometry epoch.}
  \label{fig:tut-ex15-jezero}
\end{figure}

The non-Earth body models and gravity fields are specified in
\cref{chap:dynamics}, the planetary body-fixed frames in
\cref{chap:frame-conversions}, and the positioning solution in
\cref{chap:measurements}. The underlying study is \citep{iiyama2025mars}.

% ===========================================================================
\section{Example 16: LEO PNT}
\label{sec:tut-ex16}

\textbf{Python:} \file{python/examples/ex16_leopnt.ipynb} \quad
\textbf{C++:} \file{cpp/examples/tutorials/ex16_leopnt.cc}

\paragraph{The scenario.}
The Earth companion to \cref{sec:tut-ex15}, at the same epoch
\textbf{2035-01-01}. A dense \textbf{Walker-delta 110/10/1} constellation --- 110
satellites in 10 planes --- circular at \textbf{\SI{600}{\kilo\meter}}
($a=\SI{6978}{\kilo\meter}$) and \textbf{\SI{55}{\degree}} inclination, with a
period of \SI{96.7}{\minute} (14.9 revolutions per day). Earth gravity is EGM96
truncated to $4\times4$, and the distinguishing force is Harris-Priester
atmospheric drag with a ballistic coefficient
$B = C_D A/m = \SI{0.022}{\meter\squared\per\kilogram}$ ($C_D=2.2$,
$A=\SI{1}{\meter\squared}$, $m=\SI{100}{\kilogram}$). The fleet is propagated
for two orbital periods (\SI{193}{\minute}, 194 epochs at \SI{60}{\second}).

Unlike Mars, Earth needs no special design frame: the \code{GCRF} pole sits
within a fraction of a degree of the spin axis, so an inclination specified in
\code{GCRF} is the geographic inclination.

\begin{lstlisting}[language=py, caption={Adding Harris-Priester drag to the force model.}, label={lst:tut-ex16-drag}]
# python/examples/ex16_leopnt.ipynb
CD, AREA, MASS = 2.2, 1.0, 100.0   # drag coefficient, area [m^2], mass [kg]
BSTAR = CD * AREA / MASS           # ballistic coefficient B = CD*A/m [m^2/kg]
earth = pnt.Body.Earth(N_GRAV, N_GRAV, gravity_file="EGM96.cof")

def make_dynamics(drag=True, integrator=pnt.IntegratorType.RK4, step=30.0):
    dyn = pnt.NBodyDynamics()
    dyn.set_frame(GCRF)
    dyn.set_units(pnt.SI_UNITS)
    dyn.add_body(earth)            # Earth gravity field (incl. J2)
    dyn.set_integrator(integrator)
    dyn.set_time_step(step)
    if drag:
        dyn.set_drag_coeff(BSTAR)  # Harris-Priester atmospheric drag
    return dyn
\end{lstlisting}

\paragraph{The computation.}
Four studies follow. A drag-ablation propagates one satellite twice from the
same state over five days at \SI{15}{\minute} sampling, once with and once
without drag, using the same RK8 integrator so numerical error cancels in the
difference. Coverage and PDOP are mapped on a \SI{4}{\degree} grid under a
\SI{5}{\degree} mask. A user fix is computed at San Francisco with
$\sigma=\SI{1}{\meter}$ ranging, seeded at a coarse surface point beneath the
visible set because LEO satellites sit close overhead and an Earth-centre seed
converges poorly. Finally, an onboard orbit-determination EKF over
$[\vec{r},\vec{v},b,\dot b]$ estimates one nav satellite's own state from
ionosphere-free GPS pseudoranges against a nominal Walker 24/6/2 GPS
constellation at \SI{20200}{\kilo\meter}.

\begin{lstlisting}[language=py, caption={The onboard ODTS predict step, using the native STM.}, label={lst:tut-ex16-ekf}]
# python/examples/ex16_leopnt.ipynb
dyn_filter = make_dynamics(drag=False)   # on-board model: gravity only (drag unmodeled)
SIGMA_RHO = 1.0      # iono-free pseudorange noise [m]
Q_ACC = 1e-5         # dynamics process noise (accel) [m/s^2]
Q_CLK = 5e-4         # clock-drift random walk [m/s/sqrt(s)]

xf, stm, _ = dyn_filter.propagate_stm_with_info(x[:6], TS_ODTS[k - 1], TS_ODTS[k])
x[:6] = np.asarray(xf)
x[6] += x[7] * DT_ODTS
Phi = np.eye(8)
Phi[:6, :6] = np.asarray(stm)
Phi[6, 7] = DT_ODTS
\end{lstlisting}

\paragraph{The results.}
Six figures.

Figure~1 (\cref{fig:tut-ex16-drag}) is the motivation: the position separation
between the drag and no-drag propagations, in kilometres against days. It grows
super-linearly --- drag removes energy, so the unmodelled satellite falls behind
along-track and the along-track error accumulates quadratically --- reaching
\textbf{\SI{0.74}{\kilo\meter} after five days}. For a satellite that broadcasts
its own predicted position, that is a direct, unacceptable contribution to every
user's fix, which is the reason a LEO PNT operator must model drag. The printed
budget puts it in context: at \SI{600}{\kilo\meter} the monopole is
\SI{8.20}{\meter\per\second\squared}, $J_2$ is
\SI{1.115e-2}{\meter\per\second\squared}, and drag only
\SI{2.299e-7}{\meter\per\second\squared} --- and at \SI{400}{\kilo\meter} drag
is \SI{3.277e-6}{} while at \SI{800}{\kilo\meter} it is
\SI{2.683e-8}{\meter\per\second\squared}, a factor of 100 across
\SI{400}{\kilo\meter} of altitude.

\begin{figure}[htbp]
  \centering
  \datafig[width=0.72\textwidth]{tutorials/ex16_leopnt_1.png}
  \caption{Example~16, first figure: position error accumulated by ignoring
           Harris-Priester drag at \SI{600}{\kilo\meter} over five days.}
  \label{fig:tut-ex16-drag}
\end{figure}

Figure~2 is the 110-satellite fleet in 3-D over one period in \code{GCRF},
coloured by plane over a textured Earth; figure~3 shows the Earth-fixed ground
tracks over two orbits sweeping the $\pm\SI{55}{\degree}$ band, with San
Francisco marked. Figure~4 (\cref{fig:tut-ex16-pdop}) maps median PDOP with a
contour at 6 and mean satellites in view: the characteristic footprint of an
inclined Walker constellation, a strong mid-latitude service band that thins
toward the poles. The quantitative summary over the $|\mathrm{lat}| \le
\SI{60}{\degree}$ band is a median PDOP of \textbf{6.65} and a mean of
\textbf{3.50} satellites in view --- again a fleet at the edge of sufficiency,
even at 110 satellites, because each one is visible for only a few minutes from
\SI{600}{\kilo\meter}.

\begin{figure}[htbp]
  \centering
  \datafig[width=\textwidth]{tutorials/ex16_leopnt_4.png}
  \caption{Example~16, fourth figure: median PDOP over two orbits (left,
           contour at PDOP 6) and mean satellites above the \SI{5}{\degree}
           mask (right).}
  \label{fig:tut-ex16-pdop}
\end{figure}

Figure~5 is the San Francisco fix, in the same three panels as
\cref{sec:tut-ex15}: error against time coloured by PDOP, the error CDF, and the
best-epoch sky view. San Francisco sees a minimum of 2, a mean of 4.5 and a
maximum of 6 satellites, giving a fix at \textbf{166 of 194 epochs}
(\SI{86}{\percent}); the \textbf{median 3-D error is \SI{3.3}{\meter}} with a
95th percentile of \SI{15.6}{\meter} (\SI{1.3}{\meter} horizontal,
\SI{2.9}{\meter} vertical) at a median PDOP of 4.94. The contrast with
\cref{sec:tut-ex15} is instructive: a comparable PDOP with \SI{1}{\meter}
instead of \SI{8}{\meter} ranging gives metres instead of tens of metres, so
ranging quality, not geometry, is what separates the two studies.

Figure~6 (\cref{fig:tut-ex16-odts}) is the onboard ODTS result, in three panels:
position error and filter $3\sigma$ on a log axis, clock-bias error with its
$\pm3\sigma$ band, and the number of GPS satellites in view. Looking \emph{up}
at the GPS shell from \SI{600}{\kilo\meter}, the receiver sees a minimum of 12,
a mean of 14.0 and a maximum of 16 satellites --- far better geometry than any
ground user. The filter converges from an initial error of
\textbf{\SI{2.56}{\kilo\meter}} down to a steady-state position RMS of
\textbf{\SI{0.32}{\meter}} and a clock-bias RMS of \SI{0.24}{\meter}
(\SI{0.79}{\nano\second}), with \SI{100}{\percent} of steady-state epochs inside
$3\sigma$. That is achieved with a filter whose dynamics model has \emph{no}
drag term at all: the process noise $Q_\mathrm{acc}$ absorbs it, which is the
practical answer to the problem figure~1 poses.

\begin{figure}[htbp]
  \centering
  \datafig[width=\textwidth]{tutorials/ex16_leopnt_6.png}
  \caption{Example~16, sixth figure: onboard GPS orbit determination ---
           position convergence against the filter $3\sigma$, receiver clock
           estimation, and GPS visibility from \SI{600}{\kilo\meter}.}
  \label{fig:tut-ex16-odts}
\end{figure}

The drag model, the gravity field, and the state-transition-matrix propagation
used by the filter are specified in \cref{chap:dynamics} and
\cref{chap:integration}; the EKF is in \cref{chap:filters} and the GPS
observable in \cref{chap:gnss-measurements}.

% ===========================================================================
\section{Example 17: Angles-Only OD Authored in Python}
\label{sec:tut-ex17}

\textbf{Python:} \file{python/examples/ex17_python_new_sim_example.ipynb} \quad
\textbf{C++:} none (Python only, by construction)

\paragraph{The scenario.}
The template for prototyping a new scenario without touching C++. Both the
measurement model and the estimator application are Python classes, driven by
the same C++ \py{pnt.Simulation} engine as every other example, and configured
from \file{configs/sat_bearing_odts.yaml}.

Two \cpp{Spacecraft} share the lunar environment. An \emph{observer} measures
only the bearing --- the unit line-of-sight direction --- to a \emph{target}
whose ephemeris is treated as known, and an EKF on the observer estimates the
observer's own six-state orbit $[\vec{r},\vec{v}]$; there is no clock state,
because angles are clock-independent. Measurements are taken at
\textbf{\SI{1}{\per\minute}} with a bearing noise of \textbf{2 arcseconds}, from
an a-priori of \textbf{\SI{300}{\meter}} in position and
\textbf{\SI{0.3}{\meter\per\second}} in velocity, with a process-noise
acceleration of \SI{1e-6}{\meter\per\second\squared}, over a
\textbf{\SI{6}{\hour}} arc; the notebook runs \textbf{5 Monte-Carlo trials} by
looping the seed.

Measuring a direction to a \emph{known} landmark is well posed; recovering an
\emph{unknown} target's range from bearings alone is the classic ill-posed
angles-only problem and is deliberately not attempted.

\paragraph{The computation.}
Three extension points are exercised: subclass \py{pnt.Measurement} and
implement \py{compute(self, x) -> (z, H, R)}; subclass \py{pnt.Application} with
\py{setup()} and \py{step(t)}; and register the class with
\py{pnt.register\_application} so that the \code{application:} block of the YAML
can name it. The predicted observable is evaluated through the same C++ base
path a filter uses, so the same class both generates an observation (applied to
truth, plus noise) and predicts one (applied to the filter state).

\begin{lstlisting}[language=py, caption={A measurement model written in Python.}, label={lst:tut-ex17-meas}]
# python/examples/ex17_python_new_sim_example.ipynb
class BearingMeasurement(pnt.Measurement):
    "Unit line-of-sight (bearing) to a known target position."

    def __init__(self, sigma_rad):
        pnt.Measurement.__init__(self)
        self.sigma = sigma_rad
        self.target = np.zeros(3)   # set per-epoch to the target truth position

    def compute(self, x):           # -> (z = h(x), H = dh/dx, R)
        d = self.target - x[:3]
        rn = np.linalg.norm(d)
        u = d / rn
        H = np.zeros((3, 6))
        H[:, :3] = -(np.eye(3) - np.outer(u, u)) / rn
        return u, H, self.sigma**2 * np.eye(3)
\end{lstlisting}

\begin{lstlisting}[language=py, caption={An application written in Python, then registered and run.}, label={lst:tut-ex17-app}]
# python/examples/ex17_python_new_sim_example.ipynb
class PyAnglesOdtsApp(pnt.Application):
    def __init__(self, config):
        pnt.Application.__init__(self)
        self.target = config["target"]
        self.set_frequency(config.get("frequency", 1.0 / 60))
        self.sigma = np.radians(config.get("angle_sigma_arcsec", 2.0) / 3600.0)
        self.p0 = config.get("initial_position_sigma_m", 300.0)
        self.v0 = config.get("initial_velocity_sigma_mps", 0.3)
        self.qa = config.get("process_accel_sigma_mps2", 1e-6)

    def setup(self):
        pnt.Application.setup(self)  # base schedules step() at get_frequency()
        self.dyn = make_dyn()
        self.meas = BearingMeasurement(self.sigma)

pnt.register_application("PyAnglesOdtsApp", PyAnglesOdtsApp)
sim = pnt.Simulation(cfg)            # cfg from configs/sat_bearing_odts.yaml
sim.run()
app = sim.get_agent("observer").get_application()
\end{lstlisting}

The filter force model is built from Python too --- a $8\times8$ lunar field
plus Earth and Sun, RKF45, autodiff on --- and propagated with
\py{dyn.propagate\_stm} so the covariance uses the native state-transition
matrix.

\paragraph{The results.}
Two figures.

Figure~1 (\cref{fig:tut-ex17-conv}) has two panels. The left shows the ensemble
position-error RMS over the five Monte-Carlo runs against time in hours,
descending from the \SI{300}{\meter} a-priori to the metre level as bearing
geometry accumulates; the right shows the per-axis error of one run with the
filter's own $3\sigma$ envelope shaded, clipped to $\pm\SI{400}{\meter}$, with
the estimate staying inside its covariance. The printed numbers are a
\textbf{final position error of \SI{6.31}{\meter}} RMS over the ensemble and an
\textbf{arc-wide RMS of \SI{76.84}{\meter}} --- the large arc-wide figure being
dominated by the initial transient, which is exactly what the left panel makes
visible.

\begin{figure}[htbp]
  \centering
  \datafig[width=\textwidth]{tutorials/ex17_python_new_sim_example_1.png}
  \caption{Example~17, first figure: ensemble convergence over five Monte-Carlo
           runs (left) and one run's per-axis error with the filter $3\sigma$
           bounds (right).}
  \label{fig:tut-ex17-conv}
\end{figure}

Figure~2 (\cref{fig:tut-ex17-orbit}) plots the observer's true trajectory in the
$X$--$Y$ plane of \code{MOON\_CI}, with the Moon drawn as a grey disk, coloured
by the estimate-error magnitude. It shows both the orbit geometry and where
along it the estimate is worst --- the early part of the arc --- so the reader
can connect the convergence curve to a physical location rather than only to a
time axis.

\begin{figure}[htbp]
  \centering
  \datafig[width=0.55\textwidth]{tutorials/ex17_python_new_sim_example_2.png}
  \caption{Example~17, second figure: the observer's truth orbit in
           \code{MOON\_CI}, coloured by the Python-authored EKF's position
           error.}
  \label{fig:tut-ex17-orbit}
\end{figure}

The conclusion the example is built to support is that angles-only navigation
against a \emph{known} landmark is observable and converges --- a few metres
over a \SI{6}{\hour} arc from a \SI{300}{\meter} a-priori with 2-arcsecond
bearings --- and that the whole scenario needed no C++ at all. The full
authoring contract (agents, applications, measurements, the \code{world:} block,
the event loop, and the registration mechanism in both languages) is specified
in \cref{chap:new-simulation}.

% ===========================================================================
\section{Example Index}
\label{sec:tut-summary}

\cref{tab:tut-examples} collects the seventeen examples. Python file names are
relative to \file{python/examples/} and C++ file names to
\file{cpp/examples/tutorials/}; a dash marks an example with no C++ twin (the
constellation optimiser, the Cesium and Blender tooling, and the Python-authoring
template all live in the Python layer) or with no dedicated publication.

\begin{table}[htbp]
  \centering
  \caption{The \lupnt{} tutorial examples, their files, and the publications
           behind them.}
  \label{tab:tut-examples}
  \small
  \begin{tabularx}{\textwidth}{@{}r L L L l@{}}
    \toprule
    \# & Topic & Python (\code{.ipynb}) & C++ (\code{.cc}) & Reference \\
    \midrule
    1  & Orbit propagation, frames, ELFO libration
       & \file{ex1_propagate_orbit}
       & \file{ex1_propagate_orbit}
       & \citep{iiyama2026thesis} \\
    2  & Relativistic time conversions
       & \file{ex2_time_conversions}
       & \file{ex2_time_conversions}
       & \citep{iiyama2026thesis} \\
    3  & GNSS interface: antennas, SP3 vs.\ BRDC
       & \file{ex3_gnss_interface}
       & \file{ex3_gnss_interface}
       & \citep{iiyama2024navi} \\
    4  & Plasmasphere / ionosphere delay and ray tracing
       & \file{ex4_plasmasphere}
       & \file{ex4_plasmasphere}
       & \citep{iiyama2026gcpm,iiyama2025iono} \\
    5  & GNSS measurement simulation, visibility, $C/N_0$
       & \file{ex5_gnss_measurement_sim}
       & \file{ex5_gnss_measurement_sim}
       & \citep{iiyama2024navi} \\
    6  & Sidelobe pseudorange + TDCP ODTS (UDU EKF)
       & \file{ex6_gnss_odts}
       & \file{ex6_gnss_odts}
       & \citep{iiyama2026scud,iiyama2024navi} \\
    7  & Ground-station (DSN) batch orbit determination
       & \file{ex7_groundstation_odts}
       & \file{ex7_groundstation_odts}
       & \citep{iiyama2023dcp} \\
    8  & Distributed ISL ODTS, Schmidt EKF
       & \file{ex8_isl_odts}
       & \file{ex8_isl_odts}
       & \citep{iiyama2024contact,iiyama2026edm} \\
    9  & Ephemeris and almanac fitting, bit budget
       & \file{ex9_ephemeris}
       & \file{ex9_ephemeris}
       & \citep{iiyama2025ephemeris,cortinovis2024navi} \\
    10 & Surface rover: IMU + LANS + DEM
       & \file{ex10_surface_rover}
       & \file{ex10_surface_rover}
       & \citep{iiyama2023tdcp,coimbra2025endurance} \\
    11 & Lander descent: IMU + altimeter + craters
       & \file{ex11_lander_navigation}
       & \file{ex11_lander_navigation}
       & \citep{dai2025fullstack,dai2025iongnss} \\
    12 & Constellation design, PDOP, phasing
       & \file{ex12_constellation_design}
       & ---
       & \citep{iiyama2025constellation} \\
    13 & Cesium / CZML visualisation
       & \file{ex13_cesium}
       & ---
       & --- \\
    14 & Optical navigation from horizon images
       & \file{ex14_opnav}
       & ---
       & --- \\
    15 & Mars PNT Walker constellation
       & \file{ex15_marspnt}
       & \file{ex15_marspnt}
       & \citep{iiyama2025mars} \\
    16 & LEO PNT with Harris-Priester drag
       & \file{ex16_leopnt}
       & \file{ex16_leopnt}
       & --- \\
    17 & Angles-only OD authored in Python
       & \file{ex17_python_new_sim_example}
       & ---
       & \citep{casadesus2025lupnt} \\
    \bottomrule
  \end{tabularx}
\end{table}

\cref{tab:tut-results} condenses the headline result of each example into one
line, so that a reader can find the example that already answers a question
without opening seventeen notebooks.

\begin{table}[htbp]
  \centering
  \caption{Headline result of each example, as stored in the committed notebook
           output.}
  \label{tab:tut-results}
  \small
  \begin{tabularx}{\textwidth}{@{}r L@{}}
    \toprule
    \# & Headline result \\
    \midrule
    1  & Frozen ELFO elements librate but stay bounded over \SI{720}{\dayunit};
         acceleration budget spans \SI{7.2e-1}{} (lunar monopole) to
         \SI{3.2e-8}{\meter\per\second\squared} (SRP) \\
    2  & Chebyshev $\TT-\TDB$ fit reproduces the DE440t kernel to picoseconds;
         $\LT-\TT$ drifts at $\approx\SI{56}{\micro\second\per\dayunit}$ with
         anomalistic- and synodic-month periodic terms \\
    3  & Broadcast vs.\ precise: GPS along-track RMS \SI{0.82}{\meter}, clock
         RMS \SI{1.02}{\nano\second}; Galileo \SI{0.31}{\meter} and
         \SI{0.84}{\nano\second} \\
    4  & A grazing GPS-to-lunar link carries \SI{4.97}{\TECU}, i.e.\
         \SI{0.81}{\meter} of L1 delay; ray bending contributes
         \SI{2e-6}{\meter} \\
    5  & \num{844} sidelobe links over one ELFO period, mean \num{5.3} visible
         satellites, $C/N_0$ \SIrange{20.0}{44.9}{\dBHz}, all outside the main
         beam \\
    6  & Sidelobe ODTS holds \SI{20}{}--\SI{25}{\meter} RTN RMS; TDCP improves
         RMS \SI{47.3}{} $\rightarrow$ \SI{40.1}{\meter} only if its noise is
         inflated to \SI{0.20}{\meter} \\
    7  & DSN batch converges in 5 iterations to \SI{1.64}{\meter} against a
         formal \SI{1.88}{\meter}; the SRIF smoother gives the honest
         uncertainty \\
    8  & Distributed ISL: \SI{41}{\meter} mean position (vs.\ \SI{163}{\meter}
         without exchange) and \SI{25}{\meter} clock (vs.\ \SI{410}{\meter}
         without stations); ground-only filter is 1.5--\SI{6.8}{\kilo\meter} \\
    9  & \cpp{LansEphemeris} \SI{0.35}{\meter} at \SI{428}{\bit} over
         \SI{2}{\hour}; \cpp{LansAlmanac} \SI{26.9}{\kilo\meter} at
         \SI{754}{\bit} over \SI{15}{\dayunit} \\
    10 & DEM constraint cuts the final vertical error from \SI{16.4}{} to
         \SI{1.2}{\meter} (3-D: \SI{16.7}{} to \SI{6.1}{\meter}) \\
    11 & Descent MEKF ends at \SI{5.1}{\meter}; only removing LunaNet moves it
         (to \SI{12.2}{\meter}) \\
    12 & Phasing alone moves $\mathrm{PDOP}<6$ availability from \SI{20.2}{} to
         \SI{72.0}{\percent}; N-body verification drops it to
         \SI{60.3}{\percent} \\
    13 & Three CZML scenes; the five LCRNS ELFOs check out at
         \SI{1.7}{}/\SI{17.4}{\kilo\meter} perilune/apolune altitude \\
    14 & A \SI{1.05}{px} limb-fit error is a \SI{31.9}{\kilo\meter} position
         fix; the short-arc EKF does not improve on it \\
    15 & Mars 9/3/1: median PDOP 6.74 globally, fix available
         \SI{51}{\percent} of the time at Jezero with a \SI{30.0}{\meter}
         median error \\
    16 & LEO 110/10/1: \SI{3.3}{\meter} median fix in San Francisco; onboard GPS
         EKF converges \SI{2.56}{\kilo\meter} $\rightarrow$ \SI{0.32}{\meter}
         with drag absorbed by process noise \\
    17 & Python-authored angles-only EKF converges from \SI{300}{} to
         \SI{6.3}{\meter} over \SI{6}{\hour} with 2-arcsecond bearings \\
    \bottomrule
  \end{tabularx}
\end{table}

% ===========================================================================
\section{Examples Mapped to Publications}
\label{sec:tut-publications}

\lupnt{} is the simulation backbone for the Stanford NAV Lab's lunar PNT
research, and most of the examples are reduced versions of a published study.
This section maps each group of examples onto the work that specifies its
methodology and reports its results, so that a notebook can be read alongside
the paper it came from. The simulator itself should be cited as
\citep{iiyama2023lupnt}, with the extended description in
\citep{casadesus2025lupnt}.

\begin{description}[leftmargin=2.4cm,style=nextline,font=\normalfont\itshape]
  \item[\cref{sec:tut-ex1}, \cref{sec:tut-ex2} --- fundamentals]
    The dynamics, frame, and lunicentric time-scale conventions follow
    \citep[Ch.~2]{iiyama2026thesis} and the simulator papers
    \citep{iiyama2023lupnt,casadesus2025lupnt}.
  \item[\cref{sec:tut-ex3}, \cref{sec:tut-ex5}, \cref{sec:tut-ex6} --- terrestrial GNSS]
    Lunar orbit and clock estimation from terrestrial GPS is developed in
    \citep{iiyama2024navi} and, with the stochastic-cloning UD filter used here,
    in \citep{iiyama2026scud}. See \cref{chap:filters} and
    \cref{chap:gnss-measurements}.
  \item[\cref{sec:tut-ex4} --- plasmasphere and ionosphere]
    Ionospheric and plasmaspheric delay characterisation with the Global Core
    Plasma Model is \citep{iiyama2026gcpm}, with the mitigation methodology in
    \citep{iiyama2025iono}. See \cref{chap:ionosphere}.
  \item[\cref{sec:tut-ex7}, \cref{sec:tut-ex8} --- ground-station and ISL ODTS]
    Distributed lunar ODTS and time synchronisation are
    \citep{iiyama2023dcp,iiyama2024contact}; ODTS with lunar surface stations is
    \citep{casadesus2025surface}; the inter-satellite-link autonomous clock-fault
    monitoring exercised in \cref{sec:tut-ex8} is \citep{iiyama2026edm}.
  \item[\cref{sec:tut-ex9} --- ephemeris and almanac]
    Ephemeris and almanac design for lunar navigation satellites is
    \citep{iiyama2025ephemeris}, building on the earlier parameterisation studies
    \citep{cortinovis2024navi,cortinovis2023ephemeris}. See
    \cref{chap:ephemeris-almanac}.
  \item[\cref{sec:tut-ex10}, \cref{sec:tut-ex11} --- surface and descent navigation]
    Lunar surface user positioning is \citep{iiyama2023tdcp} and
    \citep{coimbra2025endurance}; the full surface autonomy stack is
    \citep{dai2025fullstack}, with the conference version
    \citep{dai2025iongnss}.
  \item[\cref{sec:tut-ex12} --- constellation design]
    The LANS constellation trade-off analysis and staged deployment study is
    \citep{iiyama2025constellation}.
  \item[\cref{sec:tut-ex15} --- Mars PNT]
    Orbit determination and time synchronisation for a future Mars relay and
    navigation constellation is \citep{iiyama2025mars}.
  \item[General motivation]
    The overview article ``Can satellite-based radionavigation be extended to
    the Moon and other extraterrestrial bodies?'' is
    \citep{iiyama2025insidegnss}.
\end{description}

\begin{lupntnote}
\cref{sec:tut-ex13} (Cesium visualisation), \cref{sec:tut-ex14} (optical
navigation), \cref{sec:tut-ex16} (LEO PNT with Harris-Priester drag) and
\cref{sec:tut-ex17} (Python authoring template) are simulator and demonstration
examples without a dedicated publication.
\end{lupntnote}

A fuller bibliography of the group's lunar-PNT work, ordered chronologically, is
given in \cref{chap:publications}.
